\documentclass[12pt]{article}
\usepackage[utf8]{inputenc}
\usepackage[T1]{fontenc}
\usepackage{lmodern}
\usepackage{geometry}
\usepackage{hyperref}
\usepackage{microtype}
\usepackage{booktabs}
\usepackage{longtable}
\geometry{a4paper, margin=1in}
\hypersetup{colorlinks=true, linkcolor=blue, citecolor=blue, urlcolor=blue}
\title{Molecular Quantum Chemical Datasets and Databases for Machine-Learning Potentials}
\author{MLP-Xplorer}
\date{2026-05-25}

\begin{document}
\maketitle

\section*{Abstract}

The development of machine learning potentials (MLPs) capable of predicting energies, forces, and derived properties with quantum-chemical accuracy has been dramatically accelerated by the creation of large, high-quality reference datasets. This review surveys over 50 quantum-chemical databases and datasets that underpin modern MLP training and benchmarking, ranging from foundational collections of small organic molecules (e.g., QM9, QM7-X) to massive, species-specific resources (e.g., GEOM, CREMP) and reactive trajectory sets (e.g., Transition1x). We present a taxonomy that organizes these resources by molecular scope, property coverage, computational methodology, and intended use case, and discuss emerging trends such as active learning for conformational diversity, multi-fidelity reference data, and the incorporation of non-equilibrium and excited-state configurations. Critical considerations for data quality, reproducibility, and suitability for force-field development are examined, and practical guidance is offered for selecting appropriate datasets for specific MLP applications, including drug discovery, materials design, and reactive dynamics. Open challenges—including gaps in coverage of charged species, transition metals, solvation effects, and long-range interactions—are identified, and an outlook toward improved benchmarks, automated data generation, and foundation models for molecular simulation is provided. The accompanying four-column overview table (Table~\ref{tab:overview}) consolidates key details, software and methodology, and availability information for all databases covered in this review.

\section{Introduction}

The advent of machine learning potentials (MLPs) has transformed computational chemistry by enabling simulations that approach the accuracy of \textit{ab initio} methods at a fraction of the computational cost. Central to this progress is the availability of extensive, high-quality reference datasets that serve both as training corpora and as benchmarks for model validation. Over the past decade, the quantum chemistry community has produced a diverse array of such datasets, ranging from small, carefully curated collections of equilibrium structures to massive ensembles of conformers sampled from broad regions of chemical space. These resources differ not only in size and chemical coverage but also in the level of theory employed, the types of properties provided (energies, forces, dipole moments, excited-state data, etc.), and the strategies used to generate diverse molecular geometries.

This review provides a systematic survey of more than 50 quantum-chemical databases and datasets that are widely used for MLP development and benchmarking. The datasets are organized into thematic groups, including small organic molecule databases (e.g., the QM7 and QM9 families), large-scale and specialized resources (e.g., PubChemQC, Alchemy, SPICE), databases for molecular dynamics and reactive trajectories (e.g., MD17, rMD17, Transition1x, xxMD), and symmetry-aware or system-specific collections (e.g., QM-sym, COMPAS, CREMP, GEOM). A comprehensive four-column overview table (Table~\ref{tab:overview}) accompanies this text, providing for each dataset its name, a concise description, the software and methodology used, and direct download or access information.

In the following sections, we first present a taxonomy that classifies datasets according to their scope, property focus, and intended role in MLP training and evaluation. We then discuss key trends in quantum-chemical methodology observed across these datasets, including the shift from single-level to multi-fidelity reference data, the increasing inclusion of forces and non-equilibrium configurations, and the adoption of active learning for efficient sampling. Important considerations regarding data quality and reproducibility are examined, followed by practical guidance for selecting datasets suited to specific MLP development goals, such as force field parameterization for drug-like molecules, photochemical dynamics, or reactive organic chemistry. Finally, we outline open challenges and future directions, highlighting remaining gaps in coverage and promising avenues for next-generation dataset construction.

\section{Taxonomy of Dataset Types}

The datasets surveyed in this review fall into several overlapping categories based on molecular scope, the nature of the reference data, and their primary purpose in MLP development.

\subsection*{Equilibrium-Structure Databases}
These collections provide optimized geometries and associated properties for large numbers of molecules, nearly always in their ground-state equilibrium conformations. Prominent examples include the QM9~\cite{ds1-1}, QM7~\cite{ds7-32}, and QM7b~\cite{ds8-34} datasets, as well as the PubChemQC family~\cite{ds13-28,ds14-29,ds15-27}, Alchemy~\cite{ds10-6}, and PC9~\cite{ds16-26}. Such databases are essential for training models that predict molecular properties from structure, and they serve as standard benchmarks for evaluating regression accuracy. However, their restriction to a single conformation per molecule limits their direct applicability to dynamic simulations or potential energy surface (PES) learning.

\subsection*{Conformer Ensembles and Non-Equilibrium Geometries}
A second major class includes datasets that sample multiple conformers per molecule, often exploring regions of the PES away from equilibrium. The ANI family (ANI-1~\cite{ds25-7}, ANI-1x~\cite{ds26-3}, ANI-2x~\cite{ds27-8}) uses normal-mode sampling and active learning to generate millions of distorted geometries, while QM7-X~\cite{ds6-33} systematically includes up to 100 non-equilibrium conformations per equilibrium structure. Large-scale conformer databases such as GEOM~\cite{ds34-16} and CREMP~\cite{ds33-15} provide extensive ensembles for drug-like and macrocyclic molecules, respectively. These sets are particularly valuable for training MLPs that must accurately predict energies and forces over a broad region of the PES.

\subsection*{Force-Labeled Trajectory Datasets}
Datasets that explicitly include energy gradients (forces) are indispensable for developing physically accurate MLPs. The MD17 family (original, rMD17~\cite{ds20-18,ds20-19}, and CCSD(T) versions~\cite{ds20-20}) and MD22~\cite{ds22-21} provide trajectories for small to medium-sized molecules with forces computed at DFT or coupled-cluster levels. The SPICE dataset~\cite{ds12-43} offers over a million conformations of drug-like molecules with forces and multipole moments, while Transition1x~\cite{ds28-47} focuses on reactive pathways with forces in transition-state regions. The $\nabla^2$DFT dataset~\cite{ds31-23} provides forces for a substantial subset of its conformers. These resources are directly applicable to force-matching training strategies.

\subsection*{Excited-State and Spectral Databases}
A smaller but growing number of datasets target electronic excited states and optical properties. QM8~\cite{ds9-35} and bigQM7~\cite{ds17-9} provide vertical excitation energies at multiple levels of theory, and QM-symex~\cite{ds30-39} supplies detailed transition symmetries and oscillator strengths for high-symmetry molecules. The xxMD dataset~\cite{ds21-51} includes nonadiabatic dynamics trajectories with excited-state energies and forces, making it suited for MLPs targeting photochemistry.

\subsection*{Symmetry-Focused and System-Specific Sets}
Some databases are crafted to address particular chemical classes or to enforce structural constraints. The QM-sym~\cite{ds29-38} and QM-symex~\cite{ds30-39} databases ensure molecules belong to specific point groups, enabling the study of symmetry in machine learning. The COMPAS project~\cite{ds32-12,ds32-13,ds32-14} systematically covers polycyclic aromatic and heteroaromatic systems, while CREMP~\cite{ds33-15} focuses on macrocyclic peptides. These specialized resources are valuable for developing targeted MLPs and for probing transferability across chemical families.

\subsection*{Multi-Fidelity and Hierarchical Datasets}
An emerging paradigm is the provision of reference data at multiple levels of theory for the same molecular systems. Examples include QM9-G4MP2~\cite{ds2-36} (CCSD(T) quality on QM9 geometries), MultiXC-QM9~\cite{ds3-22} (76 DFT functionals plus GFN2-xTB), the ANI-1ccx subset~\cite{ds26-3} (CCSD(T) corrections), and VIB5~\cite{ds24-48} (MP2 and CCSD(T) at multiple basis sets). Such datasets enable delta-learning, transfer learning, and rigorous benchmarking of lower-level methods against higher-level references.

\section{Trends in Quantum-Chemical Methodology}

The datasets reviewed here reflect several broad methodological trends that have shaped the field of MLP development.

\subsection*{From Static to Dynamic Sampling}
Early databases such as QM7~\cite{ds7-32} and QM9~\cite{ds1-1} provided only equilibrium geometries, sufficient for property prediction but inadequate for force field training. The MD17 series~\cite{ds20-18} and later active-learning-driven sets like ANI-1x~\cite{ds26-3} and Transition1x~\cite{ds28-47} shifted the focus toward sampling the PES more broadly, including regions relevant to dynamics and reactivity. Methods such as normal-mode displacement, Wigner sampling (e.g., in WS22~\cite{ds23-50}), nudged elastic band (Transition1x~\cite{ds28-47}), and metadynamics (CREMP~\cite{ds33-15}, GEOM~\cite{ds34-16}) have become standard for generating diverse, non-equilibrium conformations.

\subsection*{Inclusion of Forces and Atomic Properties}
While energy-only databases remain common for property prediction, the inclusion of forces has become a hallmark of datasets intended for MLP training. The rMD17~\cite{ds20-19}, MD22~\cite{ds22-21}, SPICE~\cite{ds12-43}, ANI-1x~\cite{ds26-3}, and $\nabla^2$DFT~\cite{ds31-23} datasets all provide energy gradients. Additionally, datasets like AIMEl-DB~\cite{ds5-5} and Alchemy~\cite{ds10-6} supply atomic charges and multipole moments, enabling models that respect electrostatic interactions. The provision of full Hamiltonian matrices in $\nabla^2$DFT~\cite{ds31-23} opens avenues for learning electronic structure directly.

\subsection*{Multi-Fidelity Reference Data}
A notable trend is the generation of data at multiple levels of theory for the same structures, enabling hierarchical model training. The QM9-G4MP2~\cite{ds2-36} and revQM9~\cite{ds4-42} datasets upgrade the DFT-level QM9 properties to near-CCSD(T) accuracy, while MultiXC-QM9~\cite{ds3-22} provides a sweep of DFT functionals. The ANI-1ccx~\cite{ds26-3} subset and VIB5~\cite{ds24-48} combine lower-level methods with high-level coupled-cluster references. This multi-fidelity approach is crucial for developing MLPs that can extrapolate to higher accuracy with limited expensive data.

\subsection*{Scalable Conformer Generation via Active Learning}
Active learning, as exemplified by the ANI-1x and ANI-2x protocols~\cite{ds26-3,ds27-8}, has become a powerful tool for efficient data acquisition. By using ensemble disagreement as a measure of model uncertainty, subsequent rounds of quantum chemistry calculations are focused on the most informative regions of chemical space. This strategy has been adopted in the construction of Transition1x~\cite{ds28-47} and the CCSD(T) subset of ANI-1ccx~\cite{ds26-3}, maximizing the value of expensive higher-level calculations.

\subsection*{Expanding Elemental and Chemical Coverage}
Early datasets were largely restricted to the CHNO elements (e.g., ANI-1, QM9). Recent efforts have systematically added sulfur, fluorine, chlorine, and bromine (ANI-2x~\cite{ds27-8}, $\nabla^2$DFT~\cite{ds31-23}, QM-sym~\cite{ds29-38}), as well as boron and silicon (SPICE 2.0~\cite{ds12-44}). The COMPAS project~\cite{ds32-13} incorporates heteroatoms into polycyclic systems, and CREMP~\cite{ds33-15} targets macrocycles with diverse side chains. However, transition metals remain largely absent from these large-scale databases.

\subsection*{Focus on Biologically and Materially Relevant Chemical Space}
Many datasets explicitly target drug-like molecules (SPICE~\cite{ds12-43}, OrbNet Denali~\cite{ds19-25}, $\nabla^2$DFT~\cite{ds31-23}, GEOM~\cite{ds34-16}) or materials-relevant systems (COMPAS~\cite{ds32-12} for organic electronics, MD22~\cite{ds22-21} for nanostructures). This trend reflects the practical demand for MLPs that can accelerate discovery in pharmaceutical and materials science.

\section{Data Quality and Reproducibility Considerations}

The reliability of MLPs trained on quantum-chemical databases hinges critically on data quality. Several factors influence quality across the reviewed datasets.

\subsection*{Level of Theory}
The accuracy of reference calculations is a primary determinant of MLP trustworthiness. While many datasets employ density functional theory (DFT) with moderate basis sets (e.g., B3LYP/6-31G(2df,p) in QM9~\cite{ds1-1}, $\omega$B97X/6-31G(d) in ANI-1x~\cite{ds26-3}), others provide higher-quality references such as G4MP2 or CCSD(T) (QM9-G4MP2~\cite{ds2-36}, ANI-1ccx~\cite{ds26-3}, VIB5~\cite{ds24-48}). The choice of functional and basis set directly affects systematic errors such as basis set superposition error (BSSE) for non-covalent interactions, which is a known limitation of small basis sets used in the ANI family and elsewhere. Semiempirical methods (e.g., GFN2-xTB in CREMP~\cite{ds33-15} and GEOM~\cite{ds34-16}) enable massive scale but introduce larger deviations from higher-level theory.

\subsection*{Geometry Convergence and Conformation Quality}
The pipeline for generating geometries varies widely. Datasets like QM9~\cite{ds1-1} and Alchemy~\cite{ds10-6} employ multi-level optimization (e.g., PM7 followed by DFT) to avoid high-energy starting structures. Others rely on force field pre-optimization before semiempirical or DFT refinement (CREMP~\cite{ds33-15}, GEOM~\cite{ds34-16}). The bigQM7 dataset~\cite{ds17-9} uses a ConnGO workflow to preserve molecular connectivity during optimization. Convergence criteria and the presence of imaginary frequencies are infrequently reported but critical for ensuring that geometries correspond to true minima.

\subsection*{Conformational Diversity and Sampling Bias}
An inherent challenge is the balance between covering chemically relevant regions and avoiding redundancy. The ANI active learning strategy~\cite{ds26-3} addresses this by uncertainty-based selection, while the GEOM pipeline~\cite{ds34-16} uses clustering to retain conformational diversity. Nevertheless, datasets like MD17 (even revised~\cite{ds20-19}) may under-sample transition states or highly distorted structures, a gap that Transition1x~\cite{ds28-47} and WS22~\cite{ds23-50} attempt to fill. Users should be aware of the intended coverage region for each dataset.

\subsection*{Reproducibility and Open Access}
Practices for data dissemination vary. Many datasets are hosted on public repositories like Figshare, Zenodo, or GitHub (e.g., QM9~\cite{ds1-1}, ANI-1x~\cite{ds26-142}, QM7-X~\cite{ds6-81}), often including raw calculation outputs, preprocessing scripts, and standardized formats (HDF5, CSV, XYZ). Others provide dedicated websites or database backends (PubChemQC~\cite{ds13-28}). The availability of complete computational protocols (e.g., convergence thresholds, exchange-correlation functional parameters, basis set definitions) is essential for reproducing the data and for assessing its suitability for specific applications. Datasets lacking explicit descriptions of these parameters (as noted for the original MD17~\cite{ds20-18}) may lead to irreproducible benchmarks.

\section{Guidance for Selecting Datasets for MLP Development}

Given the diversity of available datasets, selecting the appropriate resource for a particular MLP project requires careful consideration of the intended application, the required property set, and the target chemical space.

\subsection*{For General-Purpose Force Fields for Small Organic Molecules}
The ANI family (ANI-1x~\cite{ds26-3}, ANI-2x~\cite{ds27-8}) remains a robust choice for training MLPs on CHNO(SFCl) molecules, providing both energies and forces with broad conformational diversity. The active learning design ensures coverage of non-equilibrium regions, and the availability of CCSD(T)-corrected subsets (ANI-1ccx~\cite{ds26-3}) enables refinement toward higher accuracy. For equilibrium property prediction, QM9~\cite{ds1-1} and its multi-fidelity extensions (QM9-G4MP2~\cite{ds2-36}, revQM9~\cite{ds4-42}) serve as standard benchmarks.

\subsection*{For Drug-Like Molecules and Non-Covalent Interactions}
The SPICE dataset~\cite{ds12-43} is specifically designed for drug-discovery relevant interactions, offering forces, multipole moments, and coverage of charged and zwitterionic species. OrbNet Denali~\cite{ds19-25} and $\nabla^2$DFT~\cite{ds31-23} provide extensive conformer ensembles for drug-like molecules, with the latter including Hamiltonian matrices for electronic structure learning. GEOM~\cite{ds34-16} offers large-scale conformer ensembles with experimental property annotations, valuable for transfer learning and thermodynamic weighting.

\subsection*{For Reactive Chemistry and Transition States}
The Transition1x dataset~\cite{ds28-47} is currently the largest dedicated resource for reactive configurations, providing forces and energies along minimum energy pathways. It is directly complementary to the ANI family and essential for training MLPs that must describe bond breaking and formation. For smaller systems, WS22~\cite{ds23-50} offers dense sampling of potential energy surfaces including barrier regions.

\subsection*{For Photochemistry and Excited-State Dynamics}
The xxMD dataset~\cite{ds21-51} is the primary resource for nonadiabatic dynamics, providing ground- and excited-state forces for photochemically active molecules. QM-symex~\cite{ds30-39} offers detailed transition symmetries for high-symmetry molecules, while QM8~\cite{ds9-35} and bigQM7~\cite{ds17-9} provide vertical excitation energies for benchmarking.

\subsection*{For Specialized Chemical Classes}
For polycyclic aromatic hydrocarbons, the COMPAS project~\cite{ds32-12} provides systematically curated data. For macrocyclic peptides, CREMP~\cite{ds33-15} covers a challenging conformational space. For symmetry-constrained systems, QM-sym~\cite{ds29-38} and QM-symex~\cite{ds30-39} offer explicit point-group labels.

\subsection*{Key Selection Criteria}
- \textit{Forces vs. energies only}: Only datasets with force labels (e.g., rMD17~\cite{ds20-19}, ANI-1x~\cite{ds26-3}, SPICE~\cite{ds12-43}, Transition1x~\cite{ds28-47}) are suitable for direct force-matching training. Energy-only datasets require alternative loss formulations.
- \textit{Chemical space}: Ensure the dataset covers the elements and molecular sizes of the target application. The ANI series covers up to eight heavy atoms (originally CHNO), while GEOM and $\nabla^2$DFT extend to drug-like molecules with 50+ atoms.
- \textit{Reference accuracy}: For benchmark validation, datasets with CCSD(T)-level references (QM9-G4MP2~\cite{ds2-36}, ANI-1ccx~\cite{ds26-3}, VIB5~\cite{ds24-48}) are recommended. For training, the cost-accuracy trade-off often favors DFT-level data with active learning.
- \textit{Conformational diversity}: Datasets that include non-equilibrium and reactive geometries (Transition1x, WS22, MD22) are necessary for MLPs intended for dynamics simulations.

\section{Open Challenges}

Despite the remarkable progress in dataset construction, several critical gaps remain.

\subsection*{Coverage of Charged and Open-Shell Species}
The majority of large-scale datasets are restricted to neutral, closed-shell molecules. Notable exceptions include SPICE~\cite{ds12-43} (which includes charged drug-like molecules) and the PubChemQC PM6 dataset~\cite{ds14-29} (which covers cationic and anionic states at the semiempirical level). However, systematic collections of radical species, multiply charged ions, and molecules with high-spin ground states are scarce, limiting the applicability of MLPs to electrochemistry, plasma chemistry, and atmospheric processes.

\subsection*{Transition Metal Chemistry}
Virtually no dataset among those reviewed provides high-quality reference data for organometallic or coordination compounds. The COMPAS project~\cite{ds32-12} touches on heterocycles with non-carbon elements but does not extend to metals. This gap severely restricts the use of MLPs in catalysis and materials science, where transition metals are ubiquitous. The development of databases covering metal–ligand bonding, spin-state energetics, and open-shell transition states remains an urgent priority.

\subsection*{Solvation and Condensed-Phase Effects}
Most datasets are limited to gas-phase molecules. While implicit solvation models (e.g., CPCM) are occasionally used, as in the BACE subset of GEOM~\cite{ds34-16}, systematic data for explicit solvent configurations or periodic systems are lacking. MLPs trained solely on gas-phase data may not transfer reliably to condensed-phase simulations without additional training or domain adaptation.

\subsection*{Long-Range Interactions and Periodic Systems}
Current datasets are dominated by small, finite molecules. Interactions such as dispersion, electrostatics at distances beyond typical molecular sizes, and periodic boundary conditions are not represented. Datasets like MD22~\cite{ds22-21} and SPICE~\cite{ds12-43} begin to address non-covalent interactions, but a comprehensive resource for intermolecular potential energy surfaces (e.g., molecular dimers, liquid-phase configurations) is needed.

\subsection*{Uncertainty Quantification and Data Consistency}
Many datasets lack explicit uncertainty estimates for their reference values. Active learning methods provide implicit uncertainty through ensemble disagreement, but systematic error bars on energies and forces are rarely reported. Additionally, mixing data from different levels of theory within a single training set (e.g., semiempirical and DFT) can introduce inconsistencies that degrade model performance if not properly treated.

\subsection*{Standardization and Interoperability}
While formats such as HDF5, XYZ, and CSV are common, there is no universal standard for metadata (e.g., charge, multiplicity, convergence thresholds, functional parameters). This makes combining multiple datasets or reproducing benchmarks challenging. Community efforts to establish common data schemas would greatly facilitate large-scale MLP training.

\section{Outlook}

The field of quantum-chemical dataset generation is evolving rapidly. Several promising directions are likely to shape the next generation of resources.

\subsection*{Foundation Models and Automated Data Curation}
Inspired by advances in natural language processing, there is growing interest in training ``foundation models'' for molecular simulations on massive, heterogeneous datasets. Automated workflows that combine conformer generation, active learning, and high-throughput calculation (as seen in the ANI pipeline~\cite{ds26-3} and $\nabla^2$DFT~\cite{ds31-23}) will continue to be refined, potentially enabling the creation of databases covering trillions of conformations across the entire organic chemical space.

\subsection*{Integration of Experimental and High-Level Theoretical Data}
Hybrid datasets that pair experimental measurements (e.g., reaction rates, binding affinities) with quantum-chemical calculations could provide more realistic benchmarks for MLPs. The GEOM dataset's inclusion of BACE experimental data~\cite{ds34-166} is a step in this direction. The increasing availability of CCSD(T)-level reference data for systematically chosen subsets (e.g., ANI-1ccx~\cite{ds26-3}) will allow MLPs to approach the ``gold standard'' of quantum chemistry at DFT cost.

\subsection*{Extension to Condensed Phases and Interfaces}
Future datasets will likely sample solvated molecules, periodic solids, and interfacial configurations. The use of implicit solvation (already present in GEOM~\cite{ds34-16}) may expand to explicit solvent shells, and periodic boundary conditions for crystals and surfaces will be essential for materials applications.

\subsection*{Inclusion of More Properties and Multi-Task Learning}
Beyond energies and forces, datasets that provide response properties (polarizabilities, hyperpolarizabilities), NMR chemical shifts, optical spectra, and transport coefficients will enable the development of multi-task MLPs capable of predicting a wide range of observables directly from atomic configurations.

\subsection*{Benchmarking and Standardized Evaluation Suites}
Robust benchmarks that test MLP performance across diverse chemical scenarios (equilibrium, reactive, solvated, excited-state) are needed to guide method development. Initiatives such as the MD17 family~\cite{ds20-18}, Transition1x~\cite{ds28-47}, and QM7-X~\cite{ds6-33} already serve this role, but a comprehensive, constantly updated benchmarking platform would accelerate progress.

\section{Conclusion}

The landscape of quantum-chemical datasets for machine learning potential development has expanded dramatically in scope, size, and sophistication. From the early QM7~\cite{ds7-32} and QM9~\cite{ds1-1} benchmarks to the billion-conformer active learning sets of the ANI family~\cite{ds25-7,ds26-3,ds27-8} and the reactive pathways of Transition1x~\cite{ds28-47}, these resources have enabled the creation of MLPs that rival or exceed the accuracy of traditional force fields for a wide range of molecular systems. The datasets reviewed here—summarized in the accompanying four-column overview table (Table~\ref{tab:overview})—represent a foundational resource for the community, but significant gaps remain, particularly in the coverage of transition metals, charged species, solvated environments, and excited-state dynamics. Addressing these gaps will require continued investment in high-throughput quantum chemistry, active learning, and data infrastructure. The convergence of automated data generation, multi-fidelity modeling, and open science practices promises to deliver increasingly capable and transferable machine learning potentials, ultimately enabling routine, accurate simulations of complex chemical processes across the physical and life sciences.

\section{Overview of Datasets}
\small
\begin{longtable}{p{0.18\textwidth}p{0.28\textwidth}p{0.28\textwidth}p{0.20\textwidth}}
\caption{Summary of datasets and databases included in the regenerated MLP-Xplorer review. The table is generated directly from the structured knowledge base.}\\
\toprule
Dataset & Details & Software and methodology & Data availability\\
\midrule
\endfirsthead
\toprule
Dataset & Details & Software and methodology & Data availability\\
\midrule
\endhead
\midrule
\multicolumn{4}{r}{Continued on next page}\\
\midrule
\endfoot
\bottomrule
\endlastfoot
QM9 \cite{ds1-1} & QM9 data set is one of the most widely used data sets, it is a collection of molecular structures and properties for 134 000 small organic molecules. QM9 targets neutral organic molecules containing up to nine non-hydrogen atoms (C, N, O, and F). & QM9 focuses on organic molecules with selection process opting for neutral molecules and excluding cations and anions. While, zwitterions are kept due to their importance in biomolecules like amino acids, heavier halogens (sulfur, bromine, chlorine, iodine) are excluded from the selection. & QM9 data set is publicly available on Figshare at https://doi.org/10.6084/m9.figshare. 978904. \cite{ds1-1} \\
QM9-G4MP2 \cite{ds2-36} & The QM9-G4MP2 database is built upon the existing QM9 data set by providing highly accurate G4MP2 (Gaussian-4 theory using reduced order perturbation) calculations for the molecular structures within QM9. The G4MP2 approach employed in this database comprises the calculations with a series of... & To generate the QM9-G4MP2 data set, the geometries of all QM9 molecules underwent reoptimization using the B3LYP/6-31G(2df,p) level theory, followed by single-point energy calculations employing CCSD(T,FC)/6-31G(d), MP2(FC)/G3MP2largeXP, RHF/mod-aug-cc-pVTZ, and RHF/mod-aug-cc-pVQZ methods. These... & The QM9-G4MP2 database is accessible in a compressed file format via Figshare at https:// doi.org/10.6084/m9.figshare.c.4351631.v1. The database includes all raw calculation outputs, scripts utilized for database construction, and pre-processed data presented... \cite{ds2-63} \\
MultiXC-QM9 \cite{ds3-22} & The MultiXC-QM9 data set significantly expands upon the well-established QM9 data set for studying small molecules. While QM9 focuses on a single functional (B3LYP) and basis set, MultiXC-QM9 offers a richer resource, encompassing 76 different DFT functionals alongside three basis sets and a... & The MultiXC-QM9 data set was constructed using xyz molecular geometries obtained from the QM9-G4MP2 data set. Molecular energies were computed using the ADF package, which implements various post-self-consistent field (SCF) methods for different functionals. & The MultiXC-QM9 data set is available on Figshare platform at https://doi.org/10.11583/ DTU.c.6185986.v3 and provides information on both molecules and reactions in separate files. Energy calculations for molecules are provided in CSV and SQLite formats. \cite{ds3-63,ds3-67} \\
revQM9 \cite{ds4-42} & The revQM9 data set is the revision of the QM9 data set with the quantum chemical properties recalculated at the aPBE0 level (which is PBE0 method reformulation with the ratio of exact exchange determined by an ML model to approximate the CCSD(T)-level energies). & The data set is the revision of the QM9 data set with the properties of ca. 130k equilibrium geometries (with up to 9 non-hydrogen atoms) calculated at a more accurate aPBE0 level with the cc-pVTZ basis set. & The revQM9 data set is available at https://zenodo.org/doi/10.5281/zenodo.10689883. The data are provided in the npz format and sample script shows how to extract the required properties from the corresponding revQM9.npz data file. \cite{ds4-71,ds4-72} \\
AIMEl-DB \cite{ds5-5} & AIMEl-DB provides atomic properties based on electronic density for 44k molecules randomly extracted from QM9 dataset, facilitating the development of ML models for direct reactivity and spectral prediction. & Initial 44 k structures (CHNO only) are extracted from QM9 with up to 9 atoms and single-point calculations are performed at B3LYP/6-31G(2df,p) level of theory with Gaussian 16. Atomic properties are generated using AIMAll package. & The dataset is publicly available at Zenodo https://zenodo.org/records/11406726. More than 30 atomic properties including energy, dipole moment, quadrupole moment and population are provided as well as molecular properties stored in Gaussian output files. \cite{ds5-60,ds5-74} \\
QM7-X \cite{ds6-33} & QM7-X is a comprehensive data set of 4.2 million structures of small organic molecules containing up to seven non-hydrogen atoms (C, N, O, S, C). It includes a systematic sampling of all stable equilibrium structures (constitutional/structural isomers and stereoisomers) and explores 100... & To generate QM7-X data set, researchers first used the GDB13 database to identify all possible structures (constitutional/structural isomers and stereoisomers) for molecules with up to seven non-hydrogen atoms. Using SMILES strings, initial 3D structures were obtained using MMFF94 force field... & The QM7-X data set is available in eight HDF5 files on Zenodo at https://doi.org/10.5281/ zenodo.4288677. Each file stores information about molecular structures using Python dictionaries. \cite{ds6-81} \\
QM7 \cite{ds7-32} & Curated from the much larger GDB-13 database containing nearly 1 billion molecules, QM7 focuses on a subset of 7165 small organic molecules. These molecules contain up to 23 atoms, with a maximum of 7 being `heavy atoms' ? & The researchers selected a subset of 7165 molecules from GDB-13 with atomization energies ranging from 800 to 2000 kcal mol-1. This subset encompasses molecules with features like double and triple bonds, cycles, and various functional groups including carboxy, cyanide, amide, alcohol, and epoxy... & The entire data set is hosted publicly on Quantum-Machine.org with data organized into several multidimensional arrays. It provides Coulomb matrices, atomization energies, atomic charge and the Cartesian coordinates of each atom in the molecules. \cite{ds7-83} \\
QM7b \cite{ds8-34} & The QM7b data set is a collection of information on over 7211 small organic molecules containing up to seven different elements: C, Cl, H, N, O, and S. It serves as an extension of the QM7 data set specifically designed for multitask learning applications in chemistry. & The creation of QM7b data set starts from the selection of 7211 small organic molecules from GDB-13 database. Initial molecular geometries were generated from SMILES strings using the universal Force Field (UFF) within the Open Babel software. & The complete QM7b data set is available as supplementary material in. It consists of two key files: coulomb.txt which contains unsorted Coulomb matrices (in atomic units) for all 7211 molecules while properties.txt stores the 14 properties associated with... \cite{ds8-34} \\
QM8 \cite{ds9-35} & The QM8 data set, introduced by Ramakrishnan et al, comprises electronic spectra data for approximately 21 786 small organic molecules. It is a carefully curated subset derived from the larger QM9 data set, which includes 134 000 molecules. & QM8 originates from the QM9 data set, containing 134 000 molecules with up to nine heavy atoms. Initially, molecules with high steric strain in their initial geometries, computed using the B3LYP/6-31G(2df,p) method, were eliminated. & The QM8 data set is accessible as a supplementary text file in. It provides indices for all molecules, facilitating the retrieval of their geometries from the QM9 data set, along with TD-DFT and CC2 excitation energies. \cite{ds9-35} \\
Alchemy \cite{ds10-6} & The Alchemy data set contains 12 quantum mechanical properties for 119 487 organic molecules with up to 14 heavy atoms, sampled from the GDB-17 database. Quantum mechanical calculations were performed using DFT (B3LYP/6-31G(2df,p)) via PySCF, providing data on molecular geometries, electronic, and... & The generation of molecular data in Alchemy was carried out using the PySCF. The quantum mechanical properties were calculated using DFT with the B3LYP functional and the 6-31G(2df,p) basis set. & The Alchemy data set is available at https://github.com/tencent-alchemy/Alchemy. \cite{ds10-94} \\
QM1B \cite{ds11-31} & QM1B, introduced by Mathiasen et al, is a colossal data set comprising one billion training examples tailored for ML applications in quantum chemistry. It encompasses molecules featuring 9-11 heavy atoms and includes properties such as energy and HOMO-LUMO gap. & The creation of the QM1B data set commenced with 1.09 million SMILES strings sourced from the GDB-11 database, focusing on molecules with 9-11 heavy atoms. Subsequently, hydrogen atoms were added to these molecules using RDKit. & The QM1B data set is accessible on the GitHub platform at https://github.com/graphcore- research/qm1b-dataset. However, it is important to highlight that QM1B is released solely for research purposes, and caution is advised for applications necessitating high... \cite{ds11-97,ds11-98} \\
SPICE \cite{ds12-43} & The SPICE data set, an acronym for Small-molecule/Protein Interaction Chemical Energies, focuses on the energetic interplay between drug-like small molecules and proteins. It is an essential resource for training models that can accurately predict forces and energies across a diverse array of... & In the computation of energies and gradients (forces) for each conformation within the SPICE data set, Psi4 serves as the primary tool. The calculations employ DFT with the B97M-D3(BJ) functional and the def2-TZVPPD basis set. & The SPICE data set is readily accessible on Zenodo at https://doi.org/10.5281/zenodo. 7338495, stored within a single HDF5 file named SPICE-1.1.2.hdf5. \cite{ds12-101} \\
PubChemQC \cite{ds13-28} & The PubChemQC database is a valuable asset for computational chemistry research, especially in material development, drug design, and ML applications. It contains electronic structures for a vast number of molecules, including approximately 3 million molecules optimized using DFT at the B3LYP/6-31G... & The construction of the PubChemQC database involved several essential steps. Firstly, data acquisition was carried out by downloading public structure data files (SDFs) from the PubChem Project FTP website, each containing information like InChI, SMILES representations, and molecular weight for... & Regarding data accessibility, the PubChemQC database was previously available at https:// pubchemqc.riken.jp as reported in. However, it has been temporarily moved to https://nakatamaho. \cite{ds13-28} \\
PubChemQC PM6 \cite{ds14-29} & The PubChemQC PM6 data sets, stands as one of the most extensive compilations in its domain, encompassing PM6 data for a staggering 221 million molecules. These data sets are based on PubChem Compounds database, and provide optimized geometries, electronic structures, and other pivotal molecular... & The construction of the PubChemQC PM6 data sets commenced with the comprehensive retrieval and parsing of the entire PubChem Compound database. Subsequently, key molecule attributes such as weight, InChI, SMILES, and formula were extracted, with exclusion criteria set for molecules exceeding 1000 g... & The PubChemQC PM6 data sets, previously accessible at http://pubchemqc.riken.jp/ pm6\_datasets.html as referenced in, are now hosted at https://nakatamaho.riken.jp/pubchemqc.riken. jp/pm6\_datasets.html for ongoing accessibility. \cite{ds14-29} \\
PubChemQC B3LYP/6-31G//PM6 \cite{ds15-27} & The PubChemQC B3LYP/6-31G//PM6 data set is a massive resource for researchers in chemistry, materials science, and drug discovery. It offers electronic properties for an astonishing 85 938 443 molecules (nearly 86 million), encompassing a broad spectrum of essential compounds and biomolecules with... & As a first step in the workflow of data generation, molecules were extracted from the PubChemQC PM6 database, ensuring a diverse set of known molecules. Open Babel software was used to generate input files compatible with the GAMESS quantum chemistry program for the target molecules. & The PubChemQC B3LYP/6-31G//PM6 data set is freely available for download at https:// nakatamaho.riken.jp/pubchemqc.riken.jp/b3lyp\_pm6\_datasets.html under a Creative Commons license. Researchers can access the data in three formats: GAMESS program input/output... \cite{ds15-105} \\
PC9 \cite{ds16-26} & The PC9 data set serves as a valuable counterpart to the widely used QM9 data set. Unlike QM9's purely theoretical approach, PC9 leverages the PubChemQC project, a vast database of 3 million real-world molecules with calculated properties at the B3LYP/6-31G(d) level of theory. & To ensure compatibility with QM9, PC9 is restricted to molecules containing a maximum of 9 `heavy atoms' (excluding H) from the elements H, C, N, O, and F. After imposing the heavy atom limit, PC9 undergoes a process to eliminate duplicates stemming from factors like enantiomers, tautomers... & The PC9 data set is readily accessible on Figshare and Zenodo. where it provides ground state and orbital properties at the B3LYP/6-31G level of theory. \cite{ds16-106,ds16-107} \\
bigQM7 \cite{ds17-1} & The bigQM7 data set is a valuable resource for researchers focused on developing ML models to predict electronic spectra of molecules. This data set includes ground-state properties and electronic spectra for over 12 880 molecules, offering a broader range of structures compared to previous data... & The creation of bigQM7 involved a detailed, multi-step process. Initially, SMILES strings for all molecules were extracted from the GDB-11 data set. & The bigQM7 data set offers multiple access points to facilitate exploration and utilization by researchers. The core data, encompassing structures, ground state properties, and electronic spectra, is conveniently available for download at... \cite{ds17-1} \\
QMug \cite{ds18-41} & The QMugs (Quantum-Mechanical Properties of Drug-like Molecules) data set provides a rich resource for researchers in drug discovery and computational chemistry. This collection curates over 665 000 molecules relevant to biology and pharmacology, extracted from the ChEMBL database. & The construction of the QMugs data set involved a multi-step process designed to capture a comprehensive picture of drug-like molecules at both the structural and QM level. The first step focused on data extraction and SMILES processing where molecules were selected from ChEMBL database based on... & The data set is hosted on the ETH Library Collection service and can be downloaded at https://doi.org/10.3929/ethz-b-000482129. The data is provided in multiple formats: (1) a summary.csv file containing computed molecular properties and annotations; (2)... \cite{ds18-114} \\
OrbNet Denali \cite{ds19-7} & The OrbNet Denali dataset is a comprehensive training collection used to develop OrbNet Denali, a ML-enhanced semiempirical method for electronic structure calculations. It includes over 2.3 million molecular geometries with corresponding energy labels calculated at the DFT and semi-empirical... & The OrbNet Denali data set is meticulously constructed to provide a broad chemical landscape relevant to both biological and organic chemistry. Initially, the ChEMBL27 database serves as a primary source, offering a wealth of chemical structures focused on SMILES strings containing elements... & The OrbNet Denali training set, comprising 2.3 million geometries and corresponding energy labels, is openly accessible on FigShare at https://doi.org/10.6084/m9.figshare.14883867. \cite{ds19-119} \\
MD17 and its later versions \cite{ds20-18} & The MD17 and its later versions serve as foundational resources for researchers benchmarking the accuracy of force fields in molecular dynamics (MD) simulations. The two versions (the 1st version (MD17) and the revised version (rMD17)) both offer a collection of MD snapshots with DFT energies and... & The original MD17 data set contains energy and force calculations derived from AIMD simulations of gas-phase molecules at room temperature. The electronic potential energies in this data set were obtained using Kohn-Sham DFT. & All data sets are available at www.sgdml.org/\#data sets. The rMD17 data set can also be accessible on Figshare at https://dx.doi.org/10.6084/m9.figshare.12672038. \cite{ds20-121,ds20-122} \\
xxMD \cite{ds21-51} & Sampling in conventional MD data sets like MD17 and rMD17 is often biased towards a narrow region of the potential energy surface (PES), primarily focusing on equilibrium structures. This limited exploration restricts the ability of neural force fields (NFFs) to accurately model chemical reactions... & The xxMD data set includes four photochemically active molecular systems: azobenzene, malonaldehyde, stilbene, and dithiophene. To investigate the complex electronic structures involved in these systems, particularly in the vicinity of deformed geometries and conical intersections, nonadiabatic... & xxMD-CASSCF and xxMD-DFT both data sets are publicly available for download on GitHub at https://github.com/zpengmei/xxMD and Zenodo at https://doi.org/10.5281/zenodo.10393859. The data is stored in compressed archives containing pre-split extended XYZ format... \cite{ds21-129} \\
MD22 \cite{ds22-21} & The MD22 data set is a comprehensive benchmark collection featuring molecules ranging from a small peptide with 42 atoms to a double-walled nanotube with 370 atoms. It comprises MD trajectories for seven systems spanning four major classes of biomolecules and supramolecular structures. & All calculations were performed using the FHI-aims electronic structure software, in conjunction with i-PI for the MD simulations. The potential energy and atomic force labels were determined at the PBE+MBD level of theory. & The MD22 data set is freely available for download at www.sgdml.org/\#data sets. \cite{ds22-130} \\
WS22 \cite{ds23-50} & WS22 is a comprehensive database focusing on ten organic molecules, varying in size and complexity, with up to 22 atoms. It includes 1.18 million geometries, encompassing both equilibrium and non-equilibrium states. & The construction of the WS22 database involves a meticulous four-step process. In the first step, researchers employed DFT to determine the equilibrium geometries for various conformations of each molecule. & The WS22 database is publicly available on Zenodo at https://doi.org/10.5281/zenodo. 7032334. \cite{ds23-133} \\
VIB5 \cite{ds24-48} & The VIB5 database is a meticulously curated collection of high-quality ab initio quantum chemical data for five small polyatomic molecules with significant astrophysical relevance: methyl chloride (CH3 Cl), methane (CH4 ), silane (SiH4 ), fluoromethane (CH3 F), and sodium hydroxide (NaOH). The... & The VIB5 database is built on a foundation of grid points that represent various nuclear configurations of the target molecules. These grid points were collected in VIB5 from the previous studies by some of the authors. & The VIB5 data is stored in a JSON file named VIB5.json, which can be downloaded from https://doi.org/10.6084/m9.figshare.16903288. The file structure organizes the data by molecule, providing information on chemical formula, number of atoms, and more. \cite{ds24-137} \\
ANI-1 \cite{ds25-7} & The ANI-1 data set provides a comprehensive collection of DFT total energy calculations of small organic molecules, encompassing approximately 20 million non-equilibrium conformations of 57 462 molecules. ANI-1 data set was used to train the general-purpose ANI-1 potential. & The creation of the ANI-1 data set involved a multi-step process including molecule selection, geometry optimization, and normal mode sampling. Initially, the GDB-11 database was used as the source, containing organic molecules with up to 11 heavy atoms (C, N, O, F). & The ANI-1 data set is stored in an HDF5 file accessible through Figshare at https://doi.org/ 10.6084/m9.figshare.5287732.v1. Additionally, a GitHub repository at https://github.com/isayev/ ANI1\_dataset provides usage details and data access instructions. \cite{ds25-139} \\
ANI-1ccx and ANI-1x \cite{ds26-3} & The ANI-1x and ANI-1ccx data sets are foundational resources used to train the universal ANI-1x and ANI-1ccx ML potentials, respectively. These data sets provide an extensive collection of millions of organic molecule conformations containing the elements carbon, hydrogen, nitrogen, and oxygen... & The ANI-1x data set was developed using an active learning procedure to enhance the diversity and accuracy of the ANI-1x potential. Initially, an ensemble of ANI models was trained using bootstrapping on a preliminary data set. & The ANI-1x and ANI-1ccx data sets are available in a single HDF5 file, accessible via Figshare at https://doi.org/10.6084/m9.figshare.c.4712477. A Python script (example\_loader.py) provided in a publicly accessible GitHub repository... \cite{ds26-3} \\
ANI-2x \cite{ds27-3} & The ANI-2x data set is used to train the ANI-2x ML model and contains 8.9 million nonequilibrium neutral singlet molecules with seven chemical elements (H, C, N, O, S, F, Cl) compared to its predecessor ANI-1x with compounds comprising only four elements (H, C, N, O), thus enabling wider ability on... & ANI-2x data set is based on GDB-11, ChEMBL, S66x8 data sets, and additional randomly-generated amino acids and dipeptides to build the data sets. Similar active learning processes as previous ANI datasets are used to generate non-equilibrium geometries with refinement on torsion, non-bonded... & The ANI-2x data sets are available on Zenodo with guidelines on use and methods used for generation https://doi.org/10.5281/zenodo.10108942. Except for the target level B97X/6-31G used to train the ANI-2x potential, additional reference data at 4 DFT levels... \cite{ds27-145} \\
Transition1x \cite{ds28-47} & The Transition1x data set is a unique resource for researchers developing ML models that can handle reactive systems. Unlike existing data sets that primarily focus on near-equilibrium configurations, Transition1x incorporates crucial data from transition regions. & The creation of Transition1x follows a meticulous workflow. First, a comprehensive database of reactant-product pairs for various organic reactions lays the foundation, providing initial geometries for reactants, products, and transition states. & The Transition1x data set is readily available on Figshare at https://doi.org/10.6084/m9. figshare.19614657.v4 and stored in a single, user-friendly HDF5 file named `Transition1x.h5'. \cite{ds28-47} \\
QM-sym \cite{ds29-38} & Symmetry plays a pivotal role in determining various molecular characteristics, such as excitation degeneracy and transition selection rules. Despite its significance, many quantum chemistry databases often overlook this crucial aspect. & QM-sym employs a two-step process, leveraging symmetry to streamline database generation. Initially, it constructs fundamental molecular frameworks based on standardized bond lengths and angles. & The QM-sym database is readily accessible through platforms like GitHub and Figshare using the links https://github.com/XI-Lab/QM-sym-database and https://doi.org/10.6084/m9.Figshare. 9638093, respectively. \cite{ds29-152,ds29-153} \\
QM-symex \cite{ds30-35} & While QM-sym focuses on larger, symmetrical organic molecules and incorporates vital symmetry information, it lacks crucial excited-state properties necessary for applications like photosensitizers and photodynamic therapy. To fill this gap, the team behind QM-sym introduced QM-symex, a database... & QM-symex inherits the rigorous data generation process established in QM-sym, ensuring the stability and symmetry of its molecules. Building upon the 135 000 structures from QM-sym, an additional 38 000 molecules were meticulously generated. & Accessible on Figshare at https://doi.org/10.6084/m9.Figshare.12815276, QM-symex provides data in XYZ format, indexed as QM\_symex\_i.xyz for ease of access, where `i' represents the order of the structure in the database. In addition to inheriting all... \cite{ds30-154} \\
2DFT \cite{ds31-23} & The ∇2DFT data set provides a comprehensive collection of approximately 16 million conformers for around 2 million drug-like molecules, featuring 8 atom types (H, C, N, O, Cl, F, Br) and up to 62 atoms. It is an extension of the original DFT data set, offering energies, forces, and various other... & The ∇2DFT data set employs a meticulous approach to generate a diverse set of conformations (1-100) for each molecule. Initially, the conformation generation method from the RDKit software suite was utilized. & The associated code and links to access various splits of the ∇2DFT data set, are freely available for download at https://github.com/AIRI-Institute/nablaDFT. \cite{ds31-23} \\
The COMPAS project \cite{ds32-12} & The COMPAS (COMputational database of Polycylic Aromatic Systems) project aims at providing structures and properties for ground-state polycyclic aromatic systems. COMPAS-1 focuses on cata-condensed polybenzenoid hydrocarbons (PBHS) with 43 k molecules, COMPAS-2 provides 0.5 million cata-condensed... & For COMPAS-1, the initial polycyclic molecules are constructed with CaGe up to 11 rings. All the geometries were optimized at GFN2-xTB level and molecules containing up to 10 rings were further optimized at B3LYP-D3(BJ)/def2-SVP with ORCA 4.2.0. & COMPAS datasets are publicly available at https://gitlab.com/porannegroup/compas where molecules and their corresponding QM properties including electronic energies, HOMO/LUMO energies and gap, ZPE, etc are hosted. An online website is developed for easier... \cite{ds32-161} \\
CREMP \cite{ds33-15} & CREMP (Conformer-rotamer ensembles of macrocyclic peptides) data set aims to explore macrocyclic peptides which is lacking in MLP field, providing 36 k representative 4-, 5-, and 6-mer homodetic cyclic peptides with up to 31.3 million unique conformers with properties calculated at semiempirical... & To generate the CREMP data set, they utilized the CREST tool to explore diverse conformational ensembles of macrocyclic peptides. They began with a set of 36 198 unique homodetic macrocyclic peptides, sampled based on key parameters such as side chains, stereochemistry, and N-methylation. & CREMP and CERMP-CycPeptMPDB data sets are available on Zenodo at https://doi.org/ 10.5281/zenodo.7931444 and https://doi.org/10.5281/zenodo.10798261. The provided data consists of CREST metadata including energies, entropies, etc and conformers with their 3D... \cite{ds33-97,ds33-165} \\
GEOM \cite{ds34-1} & The GEOM (Geometric ensemble of molecules) data set is a comprehensive resource containing high-quality conformers for 451 186 organic molecules. This includes 317 928 drug-like species with experimental data and 133 258 molecules from the QM9 data set. & The data generation process for the GEOM data set involved several key steps. It began with SMILES pre-processing, where SMILES strings were converted to their canonical forms using RDKit to ensure consistency across the dataset. & The GEOM is online at https://github.com/learningmatter-mit/geom, where both source datasets and tutorials on loading and analyzing data are provided. Conformer-level information (geometry, energy, degeneracy, etc) and species-level information (experimental... \cite{ds34-1,ds34-168} \\
tmQM \cite{ds35-46} & The tmQM data set explores transition metal-organic compound space including Werner, bioinorganic and organometallic complexes with their respective QM properties based on Cambridge structural database, enabling wide coverage of organometallic space. There are in total 86 665 complexes with 3d, 4d... & The tmQM data set was generated by extracting structures from the 2020 release of the Cambridge Structural Database (CSD) using seven filters: (1) inclusion of only mononuclear transition metal (TM) compounds, (2) containing at least one carbon and one hydrogen, (3) excluding non-metal components... & The tmQM dataset and its derivatives are available at https://github.com/bbskjelstad/ tmqm with usage information as well as timely updates. It is also part of the Quantum-Machine project and can be found at http://quantum-machine.org/datasets/. \cite{ds35-46} \\
OFF-ON dataset \cite{ds36-24} & The OFF-ON (organic fragments from organocatalysts that are non-modular) database is developed to train machine learning models aimed at predicting the behavior of structurally and conformationally diverse functional organic molecules, particularly photoswitchable organocatalysts. This dataset... & The data generation process for the OFF-ON database involved two key stages. The first stage concentrated on creating a set of molecules that represent diverse chemical environments pertinent to functional organic molecules. & The OFF-ON database is openly available at https://archive.materialscloud.org/record/ 2023.189. \cite{ds36-24,ds36-172} \\
C7O2H10-17 \cite{ds37-10} & C7O2H10-17. the C7O2H10-17 data set comprises MD trajectories for 113 randomly selected isomers of C7O2H10. & Not specified in the database record. & Available at http://quantum-machine.org/datasets/. \cite{ds37-10} \\
ISO17 \cite{ds38-17} & ISO17. The ISO17 data set extends the C7O2H10-17 data set, consisting of 129 isomers with the chemical formula C7 O2 H10 . & Not specified in the database record. & Available at http://quantum-machine.org/datasets/. \cite{ds38-17} \\
VQM24 \cite{ds39-49} & VQM24. The VQM24 data set is a comprehensive data set containing QM properties of small organic and inorganic molecules. & Not specified in the database record. & Accessible on Zenodo at https://doi.org/10.5281/zenodo.11164951. \cite{ds39-97,ds39-186} \\
QCDGE \cite{ds40-30} & QCDGE. The QCDGE database offers an extensive collection of ground and excited-state properties for 443 106 organic molecules, each containing up to ten heavy atoms, including carbon, nitrogen, oxygen, and fluorine. & Not specified in the database record. & Hosted at http://langroup.site/QCDGE/. \cite{ds40-30} \\
QM-22 \cite{ds41-40} & QM-22. The QM-22 database is a compilation of molecular data sets specifically curated for DMC calculations of the zero-point state. & Not specified in the database record. & Can be downloaded from https://github.com/jmbowma/QM-22. \cite{ds41-188} \\
CheMFi \cite{ds42-11} & CheMFi. The CheMFi data set is a multifidelity compilation of quantum chemical properties derived from a subset of the WS22 database, featuring 135 000 geometries sampled from nine distinct molecules. & Not specified in the database record. & Available at https://github.com/SM4DA/CheMFi. \cite{ds42-189} \\
QM9S \cite{ds43-37} & QM9S. The QM9S dataset, used for training and testing the DetaNet deep learning model, consists of 133 885 organic molecules derived from the QM9 dataset. & Not specified in the database record. & Hosted on Figshare at https://doi.org/10.6084/m9.figshare.24235333. \cite{ds43-190} \\
TensorMol ChemSpider \cite{ds44-45} & TensorMol ChemSpider. The TensorMol ChemSpider data set contains potential energies and forces for 3 million conformations from 15 000 different molecules composed of carbon, hydrogen, nitrogen, and oxygen. & Not specified in the database record. & Previously accessible via https://drive.google.com/drive/folders/1IfWPs7i5kfmErIRyuhGv95dSVtNFo0e\_, but noted in the paper as no longer available. \cite{ds44-191} \\
Open Molecules 2025 (OMol25) \cite{ds45-1} & The Open Molecules 2025 (OMol25) dataset is a large-scale resource for training machine learning interatomic potentials (MLIPs) for molecular chemistry, introduced by Meta FAIR. It comprises over 140 million density functional theory (DFT) single-point calculations at the ωB97M-V/def2-TZVPD level... & All calculations were performed at the ωB97M-V/def2-TZVPD level of theory using the ORCA 6.0.0 DFT package. The ωB97M-V functional is a range-separated hybrid meta-GGA, selected for its high accuracy across a broad array of quantum chemistry tasks. & The OMol25 dataset is publicly available on Hugging Face at https://huggingface.co/facebook/OMol25. Model weights are also available at the same URL. \cite{ds45-1,ds45-2} \\
SHNITSEL (Surface Hopping Nested Instances Training Set for Excited-State Learning) \cite{ds46-1} & SHNITSEL is a comprehensive data repository containing 418,870 ab-initio data points for nine organic molecules, designed to advance the development and benchmarking of machine learning (ML) models for excited-state molecular dynamics. The dataset addresses the critical lack of high-quality... & The data generation workflow combines quantum chemistry calculations and machine learning-based adaptive sampling. Quantum chemistry reference methods include: (1) CASSCF, used for all nine molecules, with active spaces ranging from CAS(2,2) for alkenes (A01-A03) to CAS(6,6) for R01 and R02, and... & The SHNITSEL dataset is available on Zenodo at https://doi.org/10.5281/zenodo.15482819. Data is organized per molecule in zip archive files, with filenames containing the molecule identifier (e.g., A01) and total number of data points. \cite{ds46-1,ds46-2} \\
ANI-1xBB \cite{ds47-1} & ANI-1xBB is a reactive machine learning interatomic potential (MLIP) dataset and model designed to improve the prediction of reaction energetics, barrier heights, and bond dissociation energies for small organic molecules. The scientific motivation is to overcome the limitations of traditional... & The data generation workflow is fully automated and consists of the following steps: (1) Initial molecules are filtered from PubChem, restricted to H, C, N, O atoms and ≤7 heavy atoms (24,001 molecules). Initial geometries are generated from SMILES using the OMEGA package. & The ANI-1xBB dataset is available for download from Carnegie Mellon University's Kilthub repository at https://kilthub.cmu.edu/articles/dataset/ANI-1xBB\_dataset/28405316. All codes used in the study are available on GitHub at... \cite{ds47-1,ds47-2} \\
QM40 \cite{ds48-1} & The QM40 dataset is a quantum mechanical (QM) dataset designed to address the scarcity of high-quality, large-scale training data for machine learning (ML) and deep learning (DL) in molecular science, particularly for drug discovery. It aims to provide a more realistic representation of drug-like... & The computational workflow for QM40 involved several steps. First, molecular SMILES strings from the ZINC database were converted into 3D structures (PDB files) using RDKit, which added hydrogen atoms and generated initial atomic coordinates for charge-neutral singlet ground states. & The QM40 dataset is publicly available on Figshare at https://doi.org/10.6084/m9.figshare.25993060.v1. It is archived in CSV file format. \cite{ds48-1,ds48-2} \\
QM9star \cite{ds49-1} & QM9star is a quantum chemical dataset of approximately 1.9 million reactive intermediates—cations, anions, and radicals—derived from the QM9 dataset. The scientific motivation is to fill a critical gap in large-scale computational data for reactive intermediates, which are key in molecular... & Data-generation workflow: Starting from equilibrium 3D structures of 134,000 neutral molecules in QM9 (B3LYP/6-31G(2df,p) level), each terminal hydrogen atom was removed while keeping the rest of the molecule unchanged. Charges and spin multiplicities were assigned: (0,2) for radicals, (1,1) for... & Repository: figshare (https://doi.org/10.6084/m9.figshare.27002905). Format: PostgreSQL database dump file. \cite{ds49-1,ds49-2} \\
MDCD20 \cite{ds50-1} & The MDCD20 dataset is a compact but chemically diverse reactive dataset of approximately 1.4 million configurations of neutral and radical molecules composed of hydrogen, carbon, nitrogen, and oxygen (H/C/N/O) elements. It was generated using an integrated molecular dynamics/coordinate... & The MDCD20 dataset was generated using an iterative MD/CD-AL framework. The initial dataset was seeded from a small subset of the Transition-1x dataset plus reactive structures from MD/CD calculations, screened by smooth overlap of atomic positions (SOAP) descriptors (n=6, l=3, s=0.5, r\_cut=5.0 Å)... & The MDCD20 dataset and MDCD-NN model are described in the research article published in CCS Chemistry (2026, Just Published, DOI: 10.31635/ccschem.026.202607339). The dataset is not explicitly deposited in a public repository in the provided text; no... \cite{ds50-1,ds50-2} \\
\end{longtable}
\normalsize



\section{Dataset and Database Catalog}

\section{Small Organic Molecule Databases: From QM7 to QM9 and Beyond}

\subsection{The QM9 Dataset: Scope and Properties}

The QM9 dataset~\cite{ds1-1} stands as one of the most widely adopted resources for benchmarking and training machine learning potentials on small organic molecules. It comprises 134~k equilibrium structures of neutral organic molecules containing up to nine non-hydrogen atoms (C, N, O, and F), drawn from the GDB-9 subset of the GDB-17 database~\cite{ds1-52}. This selection covers 621 distinct chemical formulas (stoichiometries), with C$_7$H$_{10}$O$_2$ being the most abundant (6095 constitutional isomers). Notable inclusions are biologically relevant species such as the amino acids glycine and alanine, the nucleobases cytosine, uracil, and thymine, and pharmaceutically important building blocks like pyruvic acid, piperazine, and hydroxyurea. Zwitterions are retained, but heavier halogens (S, Br, Cl, I) and charged species are excluded. The initial three-dimensional geometries were generated from SMILES strings, pre-optimized at the PM7 semi-empirical level with MOPAC~\cite{ds1-53}, and finally optimized at the B3LYP/6-31G(2df,p) level~\cite{ds1-54,ds1-55,ds1-56,ds1-57} using Gaussian~09~\cite{ds1-58}. For each molecule, the dataset provides equilibrium geometries, atomization energies, enthalpies, entropies, frontier orbital eigenvalues, dipole moments, harmonic frequencies, polarizabilities, and thermochemical energetics. For the 6095 isomers of C$_7$H$_{10}$O$_2$, energies were recalculated at the more accurate G4MP2 level. The dataset is openly accessible on Figshare~\cite{ds1-1}, with each molecule stored in an individual \texttt{dataset\_index.xyz} file.

The principal strength of QM9 lies in its size, chemical diversity, and broad adoption, which has established it as a standard benchmark for evaluating ML models. Its limitations include the use of a single density functional (B3LYP) for most properties, which introduces systematic errors relative to higher-level methods, and the restriction to equilibrium geometries, which omits the off-equilibrium configurations essential for training reactive potentials. Furthermore, the exclusion of sulfur and heavier elements limits its applicability to biologically and pharmaceutically relevant systems.

\subsection{Refinements and Extensions of QM9: QM9-G4MP2, MultiXC-QM9, revQM9, and AIMEl-DB}

Several datasets have been constructed to address the limitations of the original QM9, either by improving the accuracy of the reference calculations or by expanding the property set.

The \textbf{QM9-G4MP2} database~\cite{ds2-36} recalculates energies for all QM9 molecules using the high-level composite method G4MP2, which incorporates HF, DFT, MP2, and CCSD(T) contributions~\cite{ds2-59,ds2-61}. Geometries were re-optimized at B3LYP/6-31G(2df,p), and single-point energies were computed with CCSD(T,FC)/6-31G(d), MP2(FC)/G3MP2largeXP, RHF/mod-aug-cc-pVTZ, and RHF/mod-aug-cc-pVQZ using Gaussian~16~\cite{ds2-62}. Atomization energies were derived from separate G4MP2 calculations on isolated atoms. The database is available on Figshare~\cite{ds2-63} and includes raw outputs, preprocessing scripts, and CSV files. By providing near-CCSD(T) accuracy for the entire QM9 set, QM9-G4MP2 enables rigorous benchmarking of ML models against coupled-cluster quality targets, though the computational cost prevents its extension to larger systems.

The \textbf{MultiXC-QM9} dataset~\cite{ds3-22} takes a complementary approach by offering energies computed with 76 different DFT functionals (including PBE with SZ, DZP, and TZP basis sets) and the GFN2-xTB semi-empirical method~\cite{ds3-64}, using geometries from QM9-G4MP2~\cite{ds3-63}. Calculations were performed with the ADF package~\cite{ds3-65,ds3-66}. Beyond molecular energies, it provides reaction energies for all possible ``AB-type'' monomolecular interconversions among QM9 molecules—approximately 162 million reactions—along with bond-change information. The data are provided in CSV and SQLite formats on Figshare~\cite{ds3-67}. MultiXC-QM9 is particularly valuable for delta learning, transfer learning, and multitask learning, as it exposes relationships between low-fidelity (e.g., GFN2-xTB) and high-fidelity (e.g., PBE) methods. However, the use of fixed geometries means that the energy landscape away from equilibrium is not sampled.

The \textbf{revQM9} dataset~\cite{ds4-42} recalculates properties of approximately 130~k QM9 equilibrium geometries using the aPBE0 functional, which adjusts the fraction of exact exchange via a machine learning model to approximate CCSD(T)/cc-pVTZ atomization energies. All calculations were performed with PySCF 2.4.0~\cite{ds4-68,ds4-69,ds4-70} and the cc-pVTZ basis set; the adaptive parameter model is available on GitHub~\cite{ds4-71}. Properties include total and atomization energies, orbital energies, dipole moments, and density matrices, stored in NPZ format on Zenodo~\cite{ds4-72}. revQM9 offers systematically improved accuracy over the original DFT-level QM9 while retaining the same molecular scope, but the aPBE0 functional itself introduces a small residual error relative to CCSD(T), and the dataset does not include forces or off-equilibrium structures.

The \textbf{AIMEl-DB}~\cite{ds5-5} provides a distinct perspective by focusing on atomic properties derived from the electron density. A random subset of 44~k CHNO molecules from QM9 was used for single-point calculations at B3LYP/6-31G(2df,p) using Gaussian~16. Atomic properties (energy, dipole moment, quadrupole moment, population, etc.) were generated with AIMAll~\cite{ds5-73} and are available on Zenodo~\cite{ds5-74}. This dataset enables ML models to predict real-space descriptors of chemical reactivity and spectroscopic response directly from molecular structure, rather than only global properties. Its limitation is the modest number of molecules and the single-level DFT quality of the underlying wavefunction.

\subsection{The QM7 Family: QM7, QM7b, and QM7-X}

The earlier QM7 dataset~\cite{ds7-32} is a precursor to QM9, containing 7165 molecules selected from the GDB-13 database~\cite{ds7-77} (up to seven heavy atoms: C, N, O, S). Molecules span atomization energies from 800 to 2000~kcal~mol$^{-1}$ and include diverse functional groups and bonds. Initial geometries were generated with Open Babel~\cite{ds7-82}, and atomization energies were computed at the PBE/tier2 level with FHI-aims. The dataset is hosted on Quantum-Machine.org~\cite{ds7-83} and provides Coulomb matrices, atomic charges, and Cartesian coordinates. While QM7 is much smaller than QM9 and limited to a single property per molecule (atomization energy), its simplicity and the availability of Coulomb matrix representations made it an early benchmark for kernel-based ML models. The restriction to equilibrium structures and the moderate accuracy of the underlying PBE functional are notable limitations.

The \textbf{QM7b} dataset~\cite{ds8-34} extends QM7 to 7211 molecules and includes 14 properties per molecule, covering atomization energy (PBE0), HOMO and LUMO eigenvalues (PBE0 and GW~\cite{ds8-86}), static polarizability (PBE0 and self-consistent screening~\cite{ds8-76}), electron affinity, ionization potential, excitation energies, and maximal absorption intensity (ZINDO/s~\cite{ds8-53}). Geometries were generated from SMILES~\cite{ds8-84} via UFF~\cite{ds8-85} and optimized at the PBE level with FHI-aims. The data are available as supplementary material~\cite{ds8-88}. QM7b was designed for multitask learning, allowing models to predict multiple electronic properties simultaneously. However, like QM7, it is limited to equilibrium structures and moderate accuracy, and its small size restricts the complexity of trainable potentials.

The \textbf{QM7-X} dataset~\cite{ds6-33} represents a substantial advance in both scope and depth. It contains 4.2 million structures—all equilibrium constitutional isomers and stereoisomers of molecules with up to seven non-hydrogen atoms (C, N, O, S, Cl) from GDB-13, plus 100 non-equilibrium conformations per equilibrium structure. Non-equilibrium geometries were generated by displacing along normal mode coordinates computed at the DFTB3+MBD level~\cite{ds6-75,ds6-76}, ensuring a Boltzmann distribution of energies. Each structure is annotated with 42 properties calculated at the PBE0+MBD level using FHI-aims with tight settings~\cite{ds6-80}: energies, forces, atomic charges, dipole moments, polarizabilities, dispersion coefficients, and more. The dataset is publicly available in eight HDF5 files on Zenodo~\cite{ds6-81}. QM7-X's inclusion of off-equilibrium structures and forces makes it ideally suited for training ML potentials that accurately describe molecular dynamics and reactivity. Its limitations include a fixed level of theory (PBE0+MBD, which, while accurate for many properties, may not capture strong correlation), and the relatively small maximum molecular size (seven heavy atoms), which limits direct transferability to larger systems. Nevertheless, QM7-X remains one of the most comprehensive resources for developing and benchmarking ML potentials for small organic molecules.

\section{Large-Scale and Specialized Molecular Datasets}

\subsection{Excited-State Properties: The QM8 Dataset}

The QM8 dataset \cite{ds9-35} addresses a critical gap in quantum-chemical resources by focusing on electronic spectra rather than ground-state properties. Derived as a subset of the QM9 collection, it contains approximately 21\,786 small organic molecules with up to eight heavy atoms (C, N, O, S). QM8 provides the energies and oscillator strengths of the lowest two singlet excited states (S1 and S2), computed at both the time-dependent DFT (TD-DFT) and the second-order approximate coupled cluster (CC2) levels of theory. This dual-methodology design enables benchmarking of cheaper TD-DFT approaches against a more accurate wavefunction method, a feature particularly valuable for developing machine learning (ML) potentials that target excited-state phenomena.

The computational protocol retains the relaxed geometries from QM9 (B3LYP/6-31G(2df,p)) and performs single-point calculations with the TURBOMOLE program \cite{ds9-89}. TD-DFT calculations employ the PBE0 and CAM-B3LYP functionals with def2-SVP and def2-TZVP basis sets \cite{ds9-90,ds9-91,ds9-92}, while the RI-CC2 reference \cite{ds9-93} uses def2-TZVP. A small fraction of molecules (14) were excluded due to convergence failures or negative transition energies, underscoring the challenges of automating excited-state calculations for diverse chemical space.

For ML potential development, QM8 offers a focused test bed for predicting optical properties, but its small size and restriction to equilibrium geometries limit its utility for force-field training. The dataset is most suitable for building models that learn electronic spectra from molecular structure, and it provides a clear link to the larger QM9 family through shared molecular indices.

\subsection{Medium-Sized Organic Molecules: Alchemy}

The Alchemy dataset \cite{ds10-6} comprises 119\,487 organic molecules with up to 14 heavy atoms, sampled from the GDB-17 database. It provides 12 quantum mechanical properties computed at the B3LYP/6-31G(2df,p) level of theory using PySCF. The dataset’s scope extends beyond energy and enthalpy to include atomic charges from meta-Löwdin population analysis, which enhances the transferability of learned charge assignments across different chemical environments.

A notable strength of Alchemy is its rigorous three-stage geometry optimization: MMFF94 force field pre-optimization, HF/STO-3G refinement, and final B3LYP/6-31G(2df,p) relaxation. This pipeline reduces the risk of high-energy or unrealistic starting structures, a common pitfall in large-scale computational campaigns. The total computational cost reached approximately 3\,000\,000 CPU hours, reflecting the quality control investment.

Alchemy is well suited for training ML potentials on ground-state properties, particularly for molecules in the size range of drug-like compounds. The inclusion of atomic charges facilitates the development of models that respect electrostatics. However, the dataset does not provide forces or vibrational frequencies, which are essential for dynamics simulations. Its scope overlaps with QM9 but extends to larger molecules, making it a complementary resource for benchmarks.

\subsection{Billion-Scale Conformers: QM1B}

QM1B \cite{ds11-31} represents an extreme in dataset size, containing one billion training examples for molecules with 9–11 heavy atoms. The dataset was generated by enumerating up to 1000 conformers per molecule using the ETKDG algorithm in RDKit \cite{ds11-97} and computing HOMO-LUMO energies at the B3LYP/STO-3G level with PySCFIPU. The total conformer counts are 305.8 million, 568.7 million, and 205.4 million for 9, 10, and 11 heavy atoms, respectively.

The deliberate trade-off between accuracy and scale is evident: the minimal basis set STO-3G yields energies that are far from the chemical accuracy typically required for ML potentials. As the authors caution, QM1B is intended for research exploring the benefits of massive data volume, not for high-accuracy applications. It provides only energy and HOMO-LUMO gap, lacking forces, atomic charges, or other properties.

For ML potential development, QM1B may serve as a large-scale pretraining resource, but the low fidelity of the underlying calculations limits its direct use for force-field generation. The dataset’s value lies in probing the scaling behavior of neural network architectures and the role of data diversity versus data quality.

\subsection{Force-Focused and Drug-Like Interactions: SPICE}

The SPICE (Small-molecule/Protein Interaction Chemical Energies) dataset \cite{ds12-43,ds12-44} was designed explicitly for training models that predict forces and energies relevant to drug discovery. The initial release contained over 1.1 million conformations of drug-like molecules, dipeptides, and solvated amino acids, covering 15 elements in both charged and uncharged states. The updated version adds more than 20\,000 molecules, improves sampling of non-covalent interactions, and includes boron- and silicon-containing compounds from PubChem.

Energies and gradients are computed using the B97M-D3(BJ) functional \cite{ds12-100} with the def2-TZVPPD basis set \cite{ds12-92} in Psi4 \cite{ds12-99}. Additional properties such as dipole and quadrupole moments, atomic charges, and bond orders are provided. The inclusion of forces (energy gradients) is a crucial feature for training ML potentials that can drive molecular dynamics simulations.

SPICE’s strengths include its coverage of non-covalent interactions and its focus on biologically relevant chemical space, which directly addresses a key application domain for ML potentials. The use of a modern density functional with dispersion correction (B97M-D3(BJ)) offers a reasonable balance between accuracy and computational cost. However, the dataset is limited to single-point calculations at pre-generated conformers; it does not include reaction pathways or transition states. For ML potential development, SPICE is an excellent resource for learning energy surfaces that include weak interactions, and it is widely used for parameterizing force fields for drug-like systems.

\subsection{The PubChemQC Family: From Millions to Hundreds of Millions}

The PubChemQC project has produced a series of progressively larger databases derived from the PubChem Compound catalog. The core PubChemQC database \cite{ds13-28} contains approximately 3 million B3LYP/6-31G ground-state optimizations and over 2 million TD-DFT excited-state calculations (B3LYP/6-31+G). The computational pipeline involves multiple pre-optimization stages (PM3, HF/STO-6G) before final DFT refinement, and all molecules are treated as neutral. This dataset has been widely used for materials design and ML applications.

Expanding on this, the PubChemQC PM6 dataset \cite{ds14-29} covers 221 million molecules with PM6 calculations, including cationic, anionic, and spin-flipped states. The use of a semiempirical method (PM6) enables coverage of an enormous chemical space, but the accuracy is significantly lower than DFT. The dataset provides optimized geometries and electronic structures but no forces.

The PubChemQC B3LYP/6-31G//PM6 dataset \cite{ds15-27} further extends the scope to nearly 86 million molecules, using the PM6-optimized geometries as starting points for B3LYP/6-31G single-point calculations. For heavier elements, effective core potentials and specialized basis sets are employed. This dataset is available in multiple formats (GAMESS input/output, JSON, PostgreSQL), facilitating access by different computational tools.

Together, these databases offer a hierarchy of accuracy and coverage: from high-level DFT on a few million molecules (PubChemQC) to semiempirical treatments on hundreds of millions (PubChemQC PM6). For ML potential development, the DFT-level databases provide reliable reference data for ground-state properties, while the PM6 subsets are valuable for exploring extreme chemical diversity. However, the absence of force labels and the use of single geometries per molecule limit their direct application to force-field training. The computational cost of generating 86 million single-point calculations is substantial, yet the resulting dataset remains a unique resource for property prediction tasks.

\subsection{Bridging the Gap: PC9 as a QM9 Counterpart}

The PC9 dataset \cite{ds16-26} was created to provide a real-world complement to the purely theoretical QM9 dataset. It extracts molecules from the PubChemQC database that match QM9’s constraints (≤9 heavy atoms from {H, C, N, O, F}), yielding 99\,234 distinct compounds after removal of duplicates arising from enantiomers, tautomers, and isotopes. Remarkably, 18\,357 of these molecules overlap with QM9, while 80\,877 are unique to PC9, offering an expanded chemical space.

Unlike QM9’s closed-shell neutral molecules, PC9 includes species with higher spin multiplicities, reflecting the greater diversity of experimentally characterized compounds. The computational methodology is B3LYP/6-31G(d), providing total energies and orbital properties but not zero-point vibrational energies or thermochemical corrections, which limits direct comparability with QM9’s multiple energy terms.

For ML potential development, PC9 is valuable for testing model transferability from a small set of theoretical structures to a larger, more heterogeneous set of real molecules. The overlap subset enables controlled experiments on the effect of data distribution shifts. However, the lack of forces and the use of a single level of theory (B3LYP/6-31G(d)) mean that PC9 is best suited for evaluating models on molecular properties rather than on energy surfaces. Its primary strength lies in bridging the gap between idealized computational benchmarks and the chemical complexity found in public databases.

\subsection{Datasets for Molecular Dynamics and Excited-State Simulations}

The emergence of machine learning potentials (MLPs) capable of predicting energies and forces with near-quantum accuracy has driven the creation of diverse benchmark datasets. This section surveys a collection of datasets that span from static equilibrium properties to non-equilibrium trajectories, covering both ground and excited electronic states. We examine their scope, computational methodologies, strengths, limitations, and suitability for MLP development, highlighting comparisons where the literature supports them.

\subsubsection*{bigQM7: Electronic Spectra and Benchmarking}

The bigQM7 dataset \cite{ds17-9} extends earlier molecular datasets such as QM7 \cite{ds17-77} and QM9 \cite{ds17-1} by providing ground-state properties and electronic spectra for over 12\,880 molecules extracted from the GDB-11 database \cite{ds17-95,ds17-96}. Its primary strength lies in the systematic coverage of vertical excitation energies computed at multiple levels of theory, including ZINDO and TDDFT with basis sets ranging from 3-21G to def2-TZVP \cite{ds17-92}. A distinctive feature is the use of a three-tier connectivity-preserving geometry optimization (ConnGO) workflow \cite{ds17-108} to ensure structural integrity during high-throughput processing. Geometry optimizations were performed at the B97X-D/def2-SVP and def2-TZVP levels \cite{ds17-85,ds17-92}, and harmonic frequency analysis confirmed local minima. The dataset is accessible via a dedicated website and the NOMAD repository \cite{ds17-9}, promoting reproducibility.

For MLP development, bigQM7 is particularly valuable for training models that predict electronic properties beyond ground-state energies. The inclusion of excitation energies across multiple theoretical levels enables transfer learning and multi-fidelity modeling. However, the dataset is limited to organic molecules from the GDB-11 space (up to 7 non-hydrogen atoms) and does not include forces, restricting its direct use for reactive molecular dynamics. The reliance on density functional theory (DFT) for ground-state properties introduces typical functional and basis-set errors, which must be considered when benchmarking against higher-level methods.

\subsubsection*{QMugs and OrbNet Denali: Drug-like Chemical Space}

Two large datasets derived from the ChEMBL database provide extensive coverage of biologically relevant molecules. QMugs \cite{ds18-41} curates over 665\,000 drug-like molecules, each with three conformers optimized at the GFN2-xTB level \cite{ds18-112,ds18-113} and single-point properties computed at the B97X-D/def2-SVP level \cite{ds18-92}. The dataset includes formation energies, dipole moments, wave functions, and vibrational spectra, hosted on the ETH Library Collection \cite{ds18-114}. Its strength lies in the diversity of chemical space (3--100 heavy atoms) and the inclusion of semi-empirical and DFT data, enabling multi-level ML models. The careful filtering of ChEMBL entries for single-protein targets and neutralization of charged species ensures relevance to drug discovery.

OrbNet Denali \cite{ds19-25} extends this approach by providing over 2.3 million geometries with energy labels at both DFT and semi-empirical levels, also sourced from ChEMBL27 \cite{ds19-115,ds19-25}. It differs from QMugs in several aspects: it includes non-equilibrium configurations from normal-mode sampling and AIMD simulations at 300 K and 500 K \cite{ds19-7,ds19-112}, covers protonation states, tautomers, salt complexes, and non-covalent interactions from additional databases \cite{ds19-116,ds19-117,ds19-118}, and limits molecules to 50 atoms. The dataset is designed specifically for training the OrbNet Denali ML-enhanced semi-empirical method, with energies computed at the GFN1-xTB level \cite{ds19-112} and labels from DFT. The availability of both equilibrium and distorted geometries makes it particularly suitable for developing MLPs that must generalize to off-equilibrium configurations, a critical requirement for molecular dynamics. A potential limitation is the reliance on GFN-xTB for generating many geometries, which may introduce systematic biases compared to fully ab initio trajectories.

Compared to QMugs, OrbNet Denali offers a broader range of configurational sampling and chemical diversity (including common salts and small-molecule benchmarks), while QMugs provides richer property information such as wave functions. Both datasets are openly accessible \cite{ds18-114,ds19-119} and have been used to train robust MLPs for organic and biomolecular systems.

\subsubsection*{MD17, rMD17, and CCSD(T) Benchmark: Trajectories and Accuracy Levels}

The MD17 family \cite{ds20-18,ds20-19,ds20-20} has become a de facto benchmark for testing force field accuracy. The original MD17 dataset contains AIMD trajectories for ten small organic molecules with DFT energies and forces, but suffered from numerical noise and insufficient reporting of computational details. The revised rMD17 \cite{ds20-19} addressed these issues by recalculating properties with the PBE functional and def2-SVP basis set \cite{ds20-92,ds20-120} using stringent convergence criteria and a dense integration grid. A third version \cite{ds20-20} provides coupled-cluster energies and forces at the CCSD(T) level for a subset of molecules, using basis sets from cc-pVDZ to cc-pVTZ \cite{ds20-99}. All versions are accessible via sgdml.org and Figshare \cite{ds20-121,ds20-122}.

These datasets are invaluable for benchmarking MLP accuracy because they directly provide forces, enabling loss functions that include force matching. The contrast between DFT and CCSD(T) levels allows assessment of the intrinsic accuracy of DFT-based MLPs and the feasibility of learning at higher fidelity. However, the limited molecule count and gas-phase nature restrict generalization to larger, condensed-phase systems. The trajectory sampling is primarily near equilibrium, with the exception of the CCSD(T) set that uses the same geometries as rMD17. This narrow coverage of the potential energy surface (PES) motivated the development of datasets with broader configurational diversity.

\subsubsection*{xxMD: Nonadiabatic Dynamics and Excited States}

The extended excited-state molecular dynamics (xxMD) dataset \cite{ds21-51} was designed to overcome the equilibrium bias of MD17 by including nonadiabatic trajectories that explore regions near conical intersections and reaction pathways. It comprises four photochemically active systems: azobenzene, malonaldehyde, stilbene, and dithiophene. Two subsets are provided: xxMD-CASSCF, computed at the SA-CASSCF level of theory \cite{ds21-125} using SHARC for surface hopping \cite{ds21-124} and OpenMolcas \cite{ds21-126}, and xxMD-DFT, which recomputes the ground-state potential and forces at the M06/def2-SVP level \cite{ds21-127} using Psi4 and ASE \cite{ds21-128}. The dataset is stored in extended XYZ format and is publicly available \cite{ds21-129}.

xxMD is particularly suited for MLPs aimed at simulating photochemistry and nonadiabatic dynamics, as it provides potential energies and forces for multiple electronic states, albeit without nonadiabatic coupling vectors. The inclusion of deformed geometries and conical intersections challenges MLPs to capture strong anharmonicity and surface crossings. Its small molecule count (four systems) mirrors the limitation of MD17, but the richness of configurational space makes it a complementary benchmark. The comparison between CASSCF and DFT subsets also provides a platform for multi-fidelity learning.

\subsubsection*{MD22: Scaling to Larger Biomolecules and Nanostructures}

To address the size limitations of MD17 and xxMD, the MD22 dataset \cite{ds22-21} provides trajectories for seven systems ranging from 42 to 370 atoms, including peptides, proteins, and double-walled carbon nanotubes. Energies and forces are calculated at the PBE+MBD level of theory \cite{ds22-75,ds22-76} using FHI-aims \cite{ds22-80} and i-PI, with two basis set qualities (`light' and `tight'). The dataset is available at www.sgdml.org \cite{ds22-130}.

MD22's principal strength is its scale: it enables the training and evaluation of MLPs on systems with hundreds of atoms, testing their ability to capture long-range dispersion interactions and multi-fragment contributions. The PBE+MBD level provides a balanced treatment of van der Waals forces, which are critical for biomolecular simulations. However, the dataset covers only a single trajectory per system, limiting the diversity of conformations, and the level of theory is uniform (no multi-fidelity comparison). For MLP development, MD22 serves as an essential stress test for scaling properties.

\subsubsection*{WS22 and VIB5: High-Density Sampling and Accurate Reference Data}

The WS22 database \cite{ds23-50} offers a complementary approach by densely sampling the PES of ten organic molecules (up to 22 atoms) using Wigner sampling \cite{ds23-131} and geodesic interpolation \cite{ds23-132}. With 1.18 million geometries, it provides energies, forces, dipole moments, and electronic energies computed at the PBE0/6-311G level \cite{ds23-57}. This high density of points ensures that both equilibrium and transition-state regions are well represented, making WS22 ideal for training MLPs that require accurate potential energy surfaces across broad configurational space. The systematic generation of interpolated structures between conformers captures barriers and saddle points often missing in MD sampling. The dataset is openly accessible \cite{ds23-133}.

The VIB5 dataset \cite{ds24-48} takes a different approach by focusing on five small molecules of astrophysical importance (CH$_3$Cl, CH$_4$, SiH$_4$, CH$_3$F, NaOH) and providing theoretical best estimates (TBEs) of potential energies at the CCSD(T) complete basis set limit, along with energies and gradients at MP2/cc-pVTZ and CCSD(T)/cc-pVQZ levels \cite{ds24-134,ds24-135,ds24-136}. With over 300\,000 grid points, VIB5 offers exceptional reference quality, including core-valence correlation and diagonal Born-Oppenheimer corrections. Its multilevel nature allows for rigorous benchmarking of MLPs against near-exact quantum chemistry. The dataset is stored in a JSON file and available on Figshare \cite{ds24-137}.

Both WS22 and VIB5 are tailored for high-accuracy MLP development, but they serve different roles: WS22 provides broad configurational diversity at a moderate (DFT) level, while VIB5 offers high-fidelity reference data in specific regions of the PES. Together, they enable both the training and validation of MLPs across varying accuracy regimes.

\subsection{The ANI Family of Datasets: Scalable Coverage of Chemical Space}\label{sec:ani-family}

The ANI series of datasets represents a landmark effort in constructing large-scale quantum chemical reference data tailored specifically for training deep neural network potentials. The original ANI-1 dataset~\cite{ds25-7} provides approximately 20 million non-equilibrium conformations for 57\,462 small organic molecules containing up to eight heavy atoms (C, N, O, H). Its scope was deliberately constrained to the CHNO elements, covering neutral singlet species drawn from the GDB-11 database~\cite{ds25-95,ds25-96}. The computational methodology relied on normal-mode sampling to generate diverse geometries from optimized minima, followed by single-point energy calculations at the \(\omega\)B97X/6-31G(d) level~\cite{ds25-56} performed with Gaussian~09. A key limitation of ANI-1 is that only energies were provided; forces were not included, which restricts its direct use for training potentials that require force labels. Furthermore, the random sampling of normal modes, while efficient, may under-sample chemically relevant regions such as transition states or strongly distorted configurations. Despite these limitations, ANI-1 served as the foundation for the first general-purpose ANI potential~\cite{ds25-138} and demonstrated that deep learning models can achieve chemical accuracy across a broad domain of organic molecules.

To address the lack of force labels and improve conformational diversity, the ANI-1x dataset~\cite{ds26-3} was developed using an active learning strategy. An ensemble of preliminary ANI models was trained on an initial seed set, and subsequent rounds of sampling from databases such as GDB-11 and ChEMBL~\cite{ds26-95,ds26-96,ds26-111} were guided by the ensemble disagreement (a proxy for model uncertainty). Four complementary sampling techniques—normal mode, molecular dynamics (MD), dimer, and torsion sampling—were employed to explore diverse configurations. This iterative procedure yielded approximately 5 million conformations with associated energies and forces computed at the \(\omega\)B97X/6-31G(d) level. The resulting ANI-1x dataset thus offers a far richer representation of conformational space compared to ANI-1, particularly for non-equilibrium structures relevant for dynamics and reactivity. The ANI-1ccx dataset~\cite{ds26-3} further extended the methodology by recomputing a carefully selected 10\% subset of ANI-1x at the CCSD(T)/CBS level of theory, providing a gold-standard benchmark for coupled-cluster accuracy. The selection of these 500\,000 configurations used an active learning loop based on model uncertainty, ensuring that the high-cost calculations targeted the most informative regions of chemical space. Together, ANI-1x and ANI-1ccx enable the training of potentials that combine DFT-level coverage with CCSD(T) accuracy in critical regions~\cite{ds26-140,ds26-141}. The datasets are distributed as a single HDF5 file on Figshare~\cite{ds26-142}, and include energies at three levels (\(\omega\)B97X/6-31G, \(\omega\)B97X/def2-TZVPP, and CCSD(T)/CBS), forces, multipole moments, and charges.

The ANI-2x dataset~\cite{ds27-8} expanded the elemental scope to include sulfur, fluorine, and chlorine, thereby extending the coverage to seven elements (H, C, N, O, S, F, Cl). With 8.9 million conformations, ANI-2x incorporates additional building blocks: the S66x8 benchmark set for non-covalent interactions~\cite{ds27-144}, randomly generated amino acids and dipeptides, and specialized sampling for bulk water. The active learning protocol was refined to prioritize torsion profiles, non-bonded dimers, and water clusters, resulting in a dataset better suited for biomolecular and drug-discovery applications. All geometries were labeled at the \(\omega\)B97X/6-31G(d) level using Gaussian~09, and additional reference data at four other DFT levels are provided. Access is via Zenodo~\cite{ds27-145}. A notable limitation of the ANI family is the reliance on a relatively small basis set (6-31G(d)) for the primary DFT level, which may introduce basis-set superposition error for non-covalent interactions. Furthermore, the datasets are restricted to neutral singlet ground states, precluding charged species, open-shell systems, or electronically excited states. Nevertheless, the ANI series remains one of the most widely used resources for training universal neural network potentials due to its large scale, active learning-driven diversity, and open accessibility.

\subsection{Transition1x: Reactive Configurations for Learning Potential Energy Surfaces}\label{sec:transition1x}

Most quantum chemical datasets, including those in the ANI family, emphasize near-equilibrium or moderately distorted geometries. The Transition1x dataset~\cite{ds28-47} fills a critical gap by providing 9.6 million configurations concentrated in transition-state regions for 10\,000 organic reactions. Each data point includes energies and forces computed at the \(\omega\)B97X/6-31G(d) level using ORCA~\cite{ds28-143}, deliberately matching the level of theory used in ANI-1x~\cite{ds28-140} to enable interoperability. The core generation technique is the nudged elastic band (NEB) method~\cite{ds28-146} and its climbing-image variant (CINEB)~\cite{ds28-149}, which together produce minimum energy pathways (MEPs) connecting reactant and product geometries. Initial reactant–product pairs were sourced from a comprehensive database of organic reactions~\cite{ds28-147}. The workflow involves relaxing the end-point structures, constructing an initial path using image-dependent pair potentials (IDPP)~\cite{ds28-150}, and then relaxing the path with NEB and CINEB under the BFGS optimizer.

A key strength of Transition1x is its focus on reactive regions, which are notoriously difficult for machine learning models to capture when trained exclusively on equilibrium or randomly sampled data. The dataset thus serves both as a training resource for developing potentials that can describe bond breaking and formation, and as a benchmark for assessing the transferability of existing models to reactive dynamics. Its size (comparable to ANI-1x) and consistent methodology make it directly complementary to the ANI family. Limitations include the restriction to a single functional and basis set, and the fact that only closed-shell organic reactions are considered. Moreover, the NEB sampling may be biased toward the chosen reaction pathways, potentially underrepresenting off-pathway configurations important for full-dimensional dynamics. Nevertheless, Transition1x has already spurred progress in reactive machine learning potentials and is freely available as an HDF5 file on Figshare~\cite{ds28-151}, with dataloaders and example scripts hosted on GitLab.

\subsection{Symmetry-Aware Databases: QM-sym and QM-symex}\label{sec:qm-sym}

Symmetry is a fundamental property of molecules, influencing electronic structure, spectroscopy, and reactivity, yet many quantum chemical databases omit explicit symmetry labels. The QM-sym database~\cite{ds29-38} addresses this by providing 135\,000 structures—covering H, B, C, N, O, F, Cl, and Br—with specified point-group symmetries, primarily \(C_{2h}\), \(C_{3h}\), and \(C_{4h}\). The generation procedure employs a two-step approach: initial molecular frameworks are built using standardized bond lengths and angles consistent with the desired point group, then a genetic algorithm iteratively refines the structures to stable configurations while enforcing symmetry constraints. Each molecule is subsequently optimized at the B3LYP/6-31G(2df,p) level with Gaussian~09, and only those with successful convergence and positive vibrational frequencies are retained. This filtering ensures that the dataset contains chemically stable, symmetry-pure molecules. The QM-sym dataset provides a range of properties including point group, atomization energies, enthalpies, zero-point energies, and a series of molecular orbital energies from HOMO\(-\)5 to LUMO+5. The data are accessible via GitHub and Figshare~\cite{ds29-152,ds29-153} in XYZ format.

While QM-sym captures ground-state properties, it lacks information on excited electronic states. The companion QM-symex dataset~\cite{ds30-39} extends the library to include the first ten singlet and triplet transitions, with properties such as oscillator strength, transition energy, transition dipole moment, and symmetry labels. Building on the 135\,000 structures of QM-sym, an additional 38\,000 molecules were generated using the same symmetry-preserving protocol, again optimized at the B3LYP level with Gaussian~09. The distribution of symmetries in QM-symex is dominated by \(C_{2h}\) (46\%), followed by \(C_{3h}\) (41\%) and \(C_{4h}\) (13\%). Compared to earlier excited-state databases like QM8~\cite{ds30-35,ds30-52}, QM-symex offers explicit transition symmetries and detailed oscillator strengths, making it particularly valuable for modeling photoinduced processes and for developing machine learning models that respect symmetry constraints. Access is via Figshare~\cite{ds30-154}. A limitation of both QM-sym and QM-symex is their narrow focus on high-symmetry \(C_{nh}\) point groups; many molecules of practical interest belong to lower symmetry groups. Additionally, the computational level (B3LYP with modest basis) may restrict quantitative accuracy for excited states compared to higher-level methods like EOM-CCSD. Nevertheless, these datasets serve a unique role in enabling symmetry-aware ML models and benchmarking the role of symmetry in learning potential energy surfaces.

\subsection{Large-Scale DFT Databases: \(\nabla^2\)DFT}\label{sec:nabla2dft}

The \(\nabla^2\)DFT dataset~\cite{ds31-23} represents one of the largest publicly available collections of electronic structure data for drug-like molecules, with approximately 16 million conformers spanning about 2 million distinct molecules. The molecular source is the MOSES dataset~\cite{ds31-156}, which provides a diverse set of drug-like scaffolds and fragments. The elemental composition includes H, C, N, O, S, F, Cl, Br, resulting in up to 62 atoms per molecule. To generate conformations, RDKit was used to produce an initial pool of 1–100 conformers per molecule, followed by Butina clustering~\cite{ds31-159} to remove redundancy while retaining at least 95\% of the conformational diversity. The centroid of each cluster was selected, yielding 1–69 unique conformations per molecule. For each conformer, the Kohn–Sham DFT calculation was performed at the \(\omega\)B97X-D/def2-SVP level using the Psi4 quantum chemistry package. The dataset includes total energies, the full DFT Hamiltonian matrix, the overlap matrix, and other electronic properties. Notably, interatomic forces were computed for a subset of approximately 452\,000 molecules (2.9 million conformations), providing valuable force labels for training ML potentials. The dataset is available with open-source code and data splits at GitHub~\cite{ds31-23}.

The \(\nabla^2\)DFT dataset offers several advantages for machine learning potential development. First, its sheer size and chemical diversity—encompassing nearly 2 million drug-like molecules—provide broad coverage of the chemical space relevant for pharmaceutical applications. Second, the inclusion of the Hamiltonian matrix enables approaches that go beyond potential energy surfaces, such as learning electronic properties or constructing deep-learning density functionals. Third, the conformational sampling based on clustering reduces redundancy without sacrificing diversity. Limitations include the use of a relatively small double-zeta basis set (def2-SVP) and the \(\omega\)B97X-D functional, which, while reasonably accurate for many organic systems, may not match the accuracy of larger basis sets or higher-level wavefunction methods. Moreover, forces are available only for a fraction of the dataset, and the data are limited to ground-state, closed-shell molecules. Nonetheless, \(\nabla^2\)DFT remains a premier resource for training and benchmarking ML potentials targeting molecular dynamics and property prediction in drug-like chemical space.

\subsection{Databases for Polycyclic Aromatic Systems: The COMPAS Project}\label{sec:compas}

Polycyclic aromatic hydrocarbons (PAHs) and their heterocyclic analogues are of fundamental interest in materials science, organic electronics, and combustion chemistry. The COMPAS project (COMputational database of Polycyclic Aromatic Systems) provides a series of datasets specifically targeting these systems. COMPAS-1~\cite{ds32-12} comprises 43\,000 cata-condensed polybenzenoid hydrocarbons (PBHs) with up to 11 rings. COMPAS-2~\cite{ds32-13} expands to 500\,000 cata-condensed poly(hetero)cyclic aromatic molecules, incorporating heteroatoms, while COMPAS-3~\cite{ds32-14} focuses on 40\,000 peri-condensed PBHs. For all datasets, geometries are initially generated using the CaGe software~\cite{ds32-160} and optimized at the GFN2-xTB level. Molecules with up to 10 rings are further refined using density functional theory: COMPAS-1 and COMPAS-3 use B3LYP-D3(BJ)/def2-SVP (COMPAS-3 additionally provides CAM-B3LYP-D3BJ/aug-cc-pVDZ single points), while COMPAS-2 employs CAM-B3LYP-D3BJ/def2-SVP. All DFT calculations were performed with ORCA. The datasets include electronic energies, HOMO/LUMO energies and gaps, zero-point energies, and other properties. Data are publicly accessible via GitLab and an interactive website~\cite{ds32-161}.

The COMPAS datasets are particularly valuable for developing machine learning potentials tailored to conjugated \(\pi\)-systems, where accurate description of aromatic stabilization and non-covalent interactions is critical. The progression from PBHs to heterocycles and peri-condensed topologies allows studies of chemical trends and model transferability across increasingly complex structures. A notable strength is the provision of calculations at both semi-empirical (xTB) and DFT levels, enabling multi-fidelity learning. Limitations include the restriction to ground-state closed-shell molecules and the focus on planar or near-planar systems, which may not capture out-of-plane deformations important in some applications. Additionally, the DFT level (B3LYP with small basis) may be insufficient for quantitative predictions of charge transport or optical properties without further calibration. Nevertheless, the COMPAS project fills an important niche by providing systematically curated data for an entire class of compounds, facilitating the development of specialized ML potentials for organic electronic materials.

\subsection{Macrocyclic Peptides: The CREMP Dataset}
\label{sec:cremp}

The CREMP (Conformer-rotamer ensembles of macrocyclic peptides) dataset~\cite{ds33-15} addresses a critical gap in existing MLP training data by systematically targeting macrocyclic peptides, a molecular class underrepresented in prior quantum-chemical repositories. The primary dataset comprises 36\,198 unique homodetic 4-, 5-, and 6-mer cyclic peptides, each sampled to maximize diversity in side-chain composition, stereochemistry, and N-methylation patterns. For every peptide, the CREST workflow~\cite{ds33-164} generated extensive conformational ensembles, yielding up to 31.3 million distinct conformers overall. An auxiliary subset, CREMP-CycPeptMPDB, extends coverage to biologically relevant 6-, 7-, and 10-mer cyclic peptides curated from the CycPeptMPDB database~\cite{ds33-162}, adding 8.7 million conformers across roughly 3\,000 peptides. All properties—including energies, entropies, and three-dimensional structures—are computed at the semiempirical GFN2-xTB level~\cite{ds33-64}, consistent with the computational pipeline that begins with RDKit~\cite{ds33-97} and ETKDGv3~\cite{ds33-163} conformer generation, MMFF94~\cite{ds33-78} pre-optimization, and GFN2-xTB refinement in chloroform prior to CREST metadynamics. The data are openly released via Zenodo (doi: 10.5281/zenodo.7931444 and 10.5281/zenodo.10798261)~\cite{ds33-165}, along with usage instructions on GitHub.

The principal strength of CREMP lies in its focused coverage of macrocyclic peptide conformational space, which is notoriously challenging for both classical force fields and ab initio methods due to ring strain and extensive rotameric states. The provision of metadata from CREST (e.g., ensemble energies and entropies) enables training of MLPs that can capture thermodynamic weighting of conformers. A limitation, however, is the reliance on semiempirical GFN2-xTB reference energies; while cost-effective for generating millions of structures, the accuracy for non-covalent interactions and dispersion in large peptide rings may not match higher-level DFT. Consequently, MLPs trained on CREMP will inherit the systematic errors of the GFN2-xTB Hamiltonian, and careful validation against higher-level benchmarks (e.g., DFT-D or MP2) is advisable before application to quantitative free-energy predictions. Furthermore, the dataset does not include forces, which restricts its direct use for training reactive or gradient-based MLPs. Nonetheless, CREMP provides a unique resource for developing potentials tailored to cyclic peptide drug design, an area where conformational diversity directly impacts passive membrane permeability.

\subsection{Large-Scale Conformer Ensembles: The GEOM Dataset}
\label{sec:geom}

The GEOM (Geometric ensemble of molecules) dataset~\cite{ds34-16} offers a complementary yet broader scope, containing high-quality conformer ensembles for 451\,186 organic molecules. This collection is partitioned into drug-like species with experimental data (317\,928 molecules, sourced from the AICures initiative and the MoleculeNet benchmark~\cite{ds34-166}) and 133\,258 molecules from the QM9 dataset~\cite{ds34-1}. As with CREMP, the conformer generation pipeline employs an iterative multi-level workflow: RDKit~\cite{ds33-97} and MMFF force field provide initial geometries, the ten lowest-energy conformers are refined with GFN2-xTB, and the lowest-energy structure seeds CREST simulations~\cite{ds34-164}. A distinctive feature is the application of graph re-identification after CREST to ensure correct stereochemistry and connectivity, which is critical for drug-like molecules that may undergo rearrangement during metadynamics. For a subset of 534 molecules from the BACE dataset, further refinement via CENSO simulations~\cite{ds34-167} yields DFT-optimized geometries and single-point energies at the r2SCAN-3c/mTZVPP level (with CPCM water), providing high-accuracy references. Additionally, xTB Hessian calculations supply vibrational frequencies and thermodynamic information. The complete dataset, including conformer-level geometries, energies, degeneracies, and species-level experimental annotations, is freely available on GitHub~\cite{ds34-168}.

GEOM's primary strength is its enormous size and diversity, spanning drug-like chemical space with experimentally measured properties. The inclusion of DFT-quality energies for the BACE subset enables benchmarking of MLPs trained on the larger semiempirical set. The dataset is particularly well-suited for training MLPs that must predict relative conformer energies and populations, a key requirement for free-energy surfaces in solution. A notable limitation is the heterogeneity of the reference level: most conformers rely on GFN2-xTB, while the BACE subset uses DFT; mixing these levels in a single MLP training set may introduce inconsistencies unless careful domain separation is applied. Moreover, the absence of forces in the public release precludes direct training on potential energy surface gradients. Compared to CREMP, GEOM covers broader organic chemistry but omits macrocyclic peptides, making the two datasets complementary for covering biologically relevant cyclic structures. The drug-like focus of GEOM makes it an excellent foundation for transfer learning to peptide systems, provided that ring-specific features are accommodated.

\subsection{Transition Metal Complexes: The tmQM Dataset}
\label{sec:tmqm}

The tmQM dataset~\cite{ds35-46} systematically explores organometallic space, containing 86\,665 mononuclear transition metal complexes with 3d, 4d, and 5d metals from groups 3–12. Structures were extracted from the Cambridge Structural Database (2020 release)~\cite{ds35-169} through a series of filters (mononuclear, containing C and H, neutral or singly charged, no disorder, etc.), yielding 116\,332 candidates. After GFN2-xTB geometry optimization and additional quality filters, 86\,699 structures remained. Quantum mechanical properties were then computed at the DFT level using the TPSSh-D3BJ/def2-SVP method via Gaussian~16. The provided properties include Cartesian coordinates, electronic and dispersion energies, natural charges on the metal center, HOMO/LUMO energies and gaps, dipole moments, and polarizabilities (the latter at GFN2-xTB). Derivatives of the dataset include tmQMg (60k graphs for machine learning)~\cite{ds35-170} and tmQMg-L (30k ligand library with NBO analysis)~\cite{ds35-171}. The dataset is openly accessible on GitHub and the Quantum-Machine project website.

tmQM fills a significant void in MLP training for organometallic chemistry, where bond types (e.g., metal–ligand interactions) and electronic structure (e.g., open-shell states, variable oxidation states) pose challenges for both classical force fields and many neural network architectures. The consistent DFT level across all complexes ensures uniform accuracy, and the inclusion of metal-centered properties (natural charge, HOMO–LUMO gap) enables training of MLPs that can predict electronic features alongside energetics. A limitation is the static nature of the dataset: each complex is represented by a single optimized geometry, so conformational flexibility near the metal center is not captured. Additionally, the TPSSh/def2-SVP level, while computationally tractable, may not provide benchmark accuracy for reaction energies or spin-state splittings; higher-level corrections (e.g., coupled cluster) would be needed for quantitative catalytic modeling. For MLP development, tmQM is ideally suited for learning potential energy surfaces near equilibrium, but should be supplemented by off-equilibrium and reaction pathway data for reactive applications. Its scope—restricted to mononuclear complexes—means that multinuclear clusters and extended organometallic frameworks remain uncovered.

\subsection{Non-Equilibrium and Functional Molecules: The OFF-ON Database}
\label{sec:offon}

The OFF-ON (organic fragments from organocatalysts that are non-modular) database~\cite{ds36-24} targets structurally and conformationally diverse functional organic molecules, with a particular emphasis on photoswitchable organocatalysts and non-covalent interactions. It comprises 7\,869 equilibrium geometries (categorized as catalytic moieties, photochromic units, substituted aromatic rings, and dimers) and 67\,457 non-equilibrium geometries sampled from short molecular dynamics trajectories (5~ps each) initiated from DFTB-optimized structures. The raw MD data were reduced via farthest-point sampling to maximize diversity. Reference energies are computed using a mixed-level approach: baseline DFTB3+3ob+D3 energies are corrected to the PBE0-D3/def2-SVP level via multilinear regression, yielding a total of 75\,326 structures. The database is openly available on Materials Cloud~\cite{ds36-172}.

OFF-ON's principal strength lies in its deliberate inclusion of out-of-equilibrium conformations, which are essential for training MLPs that can accurately model reaction dynamics and conformational transitions. The focus on functional groups relevant to organocatalysis (e.g., photochromic units) provides a chemically focused resource for developing MLPs that can predict free-energy surfaces of catalytic cycles. The use of a DFTB+DFT correction scheme balances cost and accuracy; however, the correction is only as reliable as the regression model and the limited set of training points for each molecule. Furthermore, the absence of forces in the released data restricts direct force-matching approaches. Compared to GEOM and CREMP, OFF-ON is smaller in size but richer in dynamical diversity, making it complementary for training MLPs that must handle non-equilibrium processes. The dimer subset, representing non-covalent interactions, is particularly valuable for improving MLP descriptions of intermolecular forces.

\subsection{Isomer-Specific Molecular Dynamics: C7O2H10-17 and ISO17}
\label{sec:isomers}

The C7O2H10-17 dataset~\cite{ds37-10} and its extension ISO17~\cite{ds38-17} provide high-temporal-resolution molecular dynamics trajectories for isomers of the molecular formula C7O2H10. C7O2H10-17 contains trajectories for 113 randomly selected isomers, computed at 500~K with a 0.5~fs timestep using DFT (PBE) and sampled from the QM9 isomer set. ISO17 expands coverage to all 129 isomers of C7O2H10, each providing 5\,000 frames (geometries, energies, and forces) at 1~fs resolution, using PBE with Tkatchenko–Scheffler van der Waals correction in FHI-aims. Both datasets are hosted on the Quantum-Machine project website~\cite{ds37-185,ds38-185}.

These datasets are exemplary for training MLPs that require forces and energies over extended regions of conformational space, as they sample thermal fluctuations around equilibrium for each isomer. The consistent DFT level (PBE+TS) across ISO17 eliminates functional inconsistencies, and the inclusion of forces enables direct force-matching. The small stoichiometry (C7O2H10) limits chemical diversity, making these datasets best suited for benchmarking and method development rather than broad coverage. The C7O2H10-17 predecessor offers coarser sampling but at higher temporal resolution, useful for exploring ultrafast dynamics. A key limitation is the fixed temperature (500~K) and the restriction to a single molecular formula; transferability to other stoichiometries must be validated. Together, these datasets serve as standard benchmarks for evaluating MLP architectures on isomer-specific potential energy surfaces, complementing larger but force-free datasets like GEOM.

\subsection{Comprehensive Small-Molecule Quantum Properties: VQM24 and QCDGE}
\label{sec:comprehensive}

The VQM24 dataset~\cite{ds39-49} provides a broad collection of quantum-mechanical properties for 258\,242 unique constitutional isomers and 577\,705 conformers of molecules with up to five heavy atoms (C, N, O, F, Si, P, S, Cl, Br). Conformers were generated using CREST~\cite{ds39-164} with GFN2-xTB and subsequently relaxed at the B97X-D3/cc-pVDZ level in Psi4~\cite{ds39-99}. The dataset includes optimized structures, thermal properties, electronic properties, wavefunctions, and—for a subset of 10\,793 conformers—diffusion quantum Monte Carlo (DMC) energies obtained with QMCPACK~\cite{ds39-174}. VQM24 is freely accessible via Zenodo~\cite{ds39-186}.

The principal strength of VQM24 is the inclusion of DMC reference energies, which approach gold-standard accuracy for small molecules and provide a benchmark for validating lower-level DFT. The wide variety of stoichiometries (including inorganic elements) extends beyond typical organic datasets. However, the B97X-D3/cc-pVDZ level for the bulk conformers may not be sufficient for all properties, and the subset with DMC energies is relatively small. The dataset is well-suited for training MLPs on electronic properties (e.g., HOMO–LUMO gaps, dipole moments) and for benchmarking the accuracy of semiempirical or DFT-based references.

The QCDGE (Quantum Chemistry Database of Ground and Excited states)~\cite{ds40-30} focuses on ground and excited-state properties for 443\,106 organic molecules (up to ten heavy atoms: C, N, O, F) drawn from QM9, PubChemQC, and GDB-11. Ground-state optimizations and frequency calculations use B3LYP/6-31G with D3(BJ) dispersion~\cite{ds40-175}, while excited-state single-point calculations employ B97X-D/6-31G, all in Gaussian~16. Twenty-seven properties are provided, including ground-state energies, thermal corrections, and transition electric dipole moments. The database is hosted on a dedicated website~\cite{ds40-187}.

QCDGE fills a niche for excited-state MLP development, offering a large number of molecules with consistent excited-state data. The level of theory (B3LYP/6-31G for ground, B97X-D/6-31G for excited) is modest but sufficient for qualitative trends; quantitative accuracy for excitation energies is limited by the small basis set and functional choice. Compared to VQM24, QCDGE covers larger molecules (up to ten heavy atoms) but uses a lower-level theory and lacks DMC benchmarks. For MLPs targeting ground-state chemistry, QCDGE's ground-state properties are redundant with other datasets, but its inclusion of excited-state properties (e.g., transition dipole moments) enables training of models for photophysical predictions. The absence of forces and non-equilibrium geometries restricts its use to static property prediction rather than dynamics.

In summary, the datasets reviewed in this section span a wide range of chemical spaces, from macrocyclic peptides to transition metal complexes and small organic molecules. Their complementary strengths—CREMP for cyclic peptide conformations, GEOM for drug-like ensembles, tmQM for organometallics, OFF-ON for non-equilibrium structures, C7O2H10-17/ISO17 for isomer-specific dynamics, VQM24 for high-accuracy benchmarks, and QCDGE for excited states—provide a rich foundation for developing machine learning potentials. However, the heterogeneity of computational levels (semiempirical, DFT, DMC) and the frequent absence of forces in the largest datasets necessitate careful integration and validation when assembling training sets for multi-purpose MLPs.

\section{Advancing Dataset Diversity and Scale: Recent Contributions}

\subsection{High-Level Ground-State References: QM-22 and QM40}

The QM-22 dataset stands apart from typical DFT repositories by offering a compilation of molecular data expressly curated for diffusion Monte Carlo (DMC) calculations of the zero-point state~\cite{ds41-40}. Its scientific motivation is to provide benchmark-quality ground-state energies and properties for systems where high-level correlated methods are essential. Because each constituent data set within QM-22 was generated using unique computational protocols tailored to the specific molecule, the compilation does not impose a uniform quantum chemical methodology. This heterogeneity is both a strength and a limitation: it allows access to state-of-the-art DMC results across a range of small molecules, but it complicates the direct use of the entire database for training unified machine learning potentials that require consistent reference levels. The dataset is publicly available via a GitHub repository~\cite{ds41-188}, enabling reproduction and extension. For ML potential development, QM-22 is best employed as a high-accuracy validation set or as a source of specifically targeted corrections rather than as a primary training corpus, given its modest size and methodological diversity.

In contrast, QM40~\cite{ds48-1} was designed to address the scarcity of large-scale, drug-like molecular data. It comprises 162,954 neutral molecules drawn from the ZINC database, containing 10–40 atoms composed of H, C, N, O, S, F, and Cl—elements that cover 88\% of FDA-approved drug chemical space. All calculations were performed at the consistent B3LYP/6-31G(2df,p) level of theory using Gaussian16, ensuring direct comparability with QM9 and the Alchemy dataset. A unique feature is the inclusion of local vibrational mode force constants (ka) for every bond, providing a quantitative measure of bond strength that is valuable for force-field parameterization and reactive potential training. The computational workflow used GFN2-xTB pre-optimization to accelerate DFT convergence, and molecules with imaginary frequencies or unphysical force constants were removed. Properties include 16 scalar quantum mechanical quantities, optimized Cartesian coordinates, and Mulliken charges, all organized into three CSV files. The dataset is hosted on Figshare~\cite{ds48-2} and accompanied by a Python module for interactive access. Strengths include its size, drug-relevance, and consistency with established benchmarks; limitations are the restriction to neutral singlet ground states and the moderate basis set size, which may not capture dispersion or anionic effects. For ML potential development, QM40 is ideally suited for training models that target drug-like chemical space, and its smooth concatenation with QM9 enables seamless multi-scale training.

\subsection{Excited-State and Multi-Fidelity Benchmarks: CheMFi, QM9S, and SHNITSEL}

Excited-state phenomena introduce additional complexity to ML potential development, as models must capture multiple electronic states, nonadiabatic couplings, and spin-orbit interactions. Three datasets have recently emerged to fill this gap, each with distinct design philosophies.

CheMFi~\cite{ds42-11} is a multi-fidelity collection focused on vertical excitation energies, oscillator strengths, dipole moments, and ground-state energies. It samples 135,000 geometries from nine molecules drawn from the WS22 database and provides five fidelity tiers defined by increasing basis set size (STO-3G through def2-TZVP), all computed at the TD-DFT/CAM-B3LYP level with ORCA. The inclusion of computational times for each fidelity level enables rigorous benchmarking of multi-fidelity models, which can learn to combine cheap low-fidelity data with expensive high-fidelity calculations. CheMFi is publicly available on GitHub~\cite{ds42-189}. Its primary strength is the systematic coverage of fidelity, allowing ML architectures to exploit correlation across levels while reducing overall cost. A limitation is the restricted molecular diversity (only nine molecules), which limits transferability to unseen chemical classes.

QM9S~\cite{ds43-37} extends the QM9 molecular set (133,885 molecules) beyond scalar properties to include vectors (electric dipole), second-order tensors (Hessian, quadrupole, polarizability), and third-order tensors (octupole, first hyperpolarizability), all recalculated at the B3LYP/def-TZVP level with Gaussian16. It also provides IR, Raman, and UV-Vis spectra via frequency analysis and TD-DFT. This rich tensorial information is crucial for developing ML potentials that predict response properties and for training models that respect rotational equivariance. The dataset is hosted on Figshare~\cite{ds43-190}. Its strengths lie in the breadth of properties and the large, consistent chemical space; the limitation is that geometries are re-optimised rather than sampled off-equilibrium, so the data do not cover reactive or non-equilibrium regions required for dynamics.

SHNITSEL~\cite{ds46-1} is a comprehensive repository of 418,870 ab-initio data points for nine organic molecules chosen to represent diverse photochemical behaviours—photoisomerization, ring-opening, ultrafast internal conversion, and spin-orbit coupling. It combines static (sampled) and dynamic (trajectory) geometries, providing energies, forces, dipole moments, transition dipoles, nonadiabatic couplings (NACs), and spin-orbit couplings (SOCs) for multiple singlet and triplet states. The majority of data (97\%) comes from multi-reference methods including CASSCF, MR-CISD, CASPT2, and ADC(2), with adaptive sampling guided by SchNarc models to efficiently explore undersized regions. Dynamic trajectories were generated using the SHARC program at the SA(3)-CASSCF level. The dataset is available on Zenodo~\cite{ds46-2}. Its strengths include the high level of theory, the inclusion of nonadiabatic and spin-orbit couplings, and the standardized format (NetCDF4) for ease of use. Limitations stem from the small molecule set and the methodological heterogeneity (different active spaces and basis sets for different molecules), which challenge transferable model training. SHNITSEL sets a new benchmark for excited-state ML, providing the necessary labels for developing models that can drive nonadiabatic dynamics.

\subsection{Large-Scale Multi-Purpose DFT Datasets: OMol25 and TensorMol ChemSpider}

Two datasets represent efforts to provide large-scale DFT reference data for general-purpose ML potentials. TensorMol ChemSpider~\cite{ds44-45} introduced 3 million conformations from 15,000 CHNO molecules, with energies and forces computed at B97X-D/6-311G level and geometries sampled via metadynamics. While pioneering in its size, the dataset was previously accessible via a Google Drive link now reported as no longer available~\cite{ds44-191}. This regrettable loss highlights the importance of persistent data archiving. The molecular scope—limited to four elements—and the single functional/basis set represent both a strength (consistent methodology) and a limitation (limited elemental diversity).

OMol25~\cite{ds45-1} by Meta FAIR marks a step change in scale and diversity. With over 140 million single-point DFT calculations at ωB97M-V/def2-TZVPD, it covers 83 elements (the entire first 83 of the periodic table) and approximately 83 million unique systems ranging from small molecules and biomolecules to metal complexes and electrolytes, up to 350 atoms. The methodology is uniform and highly modern: ωB97M-V is a range-separated hybrid meta-GGA selected for broad accuracy; def2-TZVPD with ECPs accommodates heavy elements; RI-J and COSX acceleration with tight convergence ensures consistency. Unrestricted KS formalism was applied to all non-singlet and transition-metal systems, with a deliberate HOMO–LUMO rotation to break spin symmetry. Data generation employed at least 14 distinct sampling strategies across four domains, including classical MD, RPMD, Architector-generated metal complexes, AFIR reactive sampling, and recomputation of existing community datasets (ANI-2X, Transition-1X, SPICE2, etc.) at the same level. This recomputation is a critical feature, enabling direct comparison of models without the confound of differing computational protocols. The dataset is released under CC BY 4.0 on Hugging Face~\cite{ds45-6}, with model weights under a commercially permissive license. The massive computational cost (6.6 billion CPU core-hours) makes independent reproduction infeasible, underscoring the importance of permanent public hosting. Strengths include unprecedented elemental and structural diversity, consistent high-quality DFT, and the inclusion of partial charges, spin densities, orbital energies, Fock matrices, and densities—beyond energies and forces. Limitations arise from the use of a single functional (no multi-level or error estimates) and the fact that the data are largely single-point calculations; dynamic trajectories are not included. For ML potential development, OMol25 is poised to become a foundational training resource, much as image datasets like ImageNet transformed computer vision, given its breadth and rigorous quality control.

\subsection{Reactive Machine Learning Potentials: ANI-1xBB}

ANI-1xBB~\cite{ds47-1} addresses a well-recognized deficiency in existing ML potential datasets: the near-absence of non-equilibrium and reactive geometries. It introduces an automated bond-breaking workflow that elongates each unique bond in 24,001 small organic molecules (H, C, N, O, ≤7 heavy atoms) to three times its equilibrium length over 15 steps, with geometry optimization and fixed-distance molecular dynamics at each step. This procedure generated 13,144,877 geometries, each labeled with energies, forces, dipole moments, and fractional orbital densities (NFOD) at finite electronic temperatures (0, 1000, and 5000 K) using FT-DFT with the B97-3c functional in ORCA. The finite-temperature treatment is crucial to handle open-shell character during bond dissociation. The sampling workflow includes a minimax algorithm to select diverse MD snapshots and discards non-convergent calculations. The dataset is available on Kilthub~\cite{ds47-6}, and the workflow code on GitHub~\cite{ds47-1}. When combined with equilibrium data (e.g., ANI-1x), models trained on ANI-1xBB show significant improvements in predicting reaction barrier heights, bond dissociation energies, and pericyclic reaction pathways. Strengths include the automated generation of reactive conformers, the use of a consistent quantum chemical method, and the integration with the established ANI model family. Limitations are the restricted elemental scope and molecular size, and the reliance on a single, moderately cheap DFT functional (B97-3c) that may not capture multi-reference character accurately. For ML potential development, ANI-1xBB provides a template for generating reaction-aware training data and demonstrates that explicitly including bond-breaking structures substantially enhances model transferability to reaction energetics.

\subsection{QM9star: Reactive Intermediates for Charged and Radical Species}

The QM9star dataset~\cite{ds49-1,ds49-2,ds49-3,ds49-4} represents a purpose-built extension of the widely used QM9 collection, targeting a critical gap in quantum chemical resources for reactive intermediates. Whereas most existing datasets focus on neutral, closed-shell molecules, QM9star provides approximately 1.9 million structures spanning cations, anions, and radicals, derived by systematic removal of terminal hydrogen atoms from the 134,000 equilibrium neutral molecules in QM9. The final dataset comprises 436,000 cations, 721,000 anions, 731,000 radicals, and 120,000 re-optimized neutral molecules, all with up to nine heavy atoms (C, N, O, F). Each entry includes 39 properties: global quantities such as orbital energies, vibrational frequencies, thermodynamic corrections, dipole moments, and polarizability, as well as local descriptors including Mulliken and NPA charges, spin densities, atomic forces, and NBO bond orders.

The computational methodology employed Gaussian~16 at the B3LYP-D3(BJ)/6-311+G(d,p) level of DFT, followed by NBO analysis on all successfully optimized structures. Quality control was stringent: structures failing geometry convergence, exhibiting imaginary frequencies, or with maximum forces exceeding 0.001~Eh/bohr were discarded. The resulting dataset contains 1,915,870 topologically distinct species and 2,008,806 distinct 3D conformations. A DimeNet++ model trained on formation energies achieved an MAE of 0.235~kcal/mol, approaching chemical accuracy, demonstrating the suitability of QM9star for ML potential development across multiple charge states.

The dataset is openly accessible via figshare under a Creative Commons CC BY-NC-ND~4.0 license, distributed as a PostgreSQL database dump. The accompanying code repository provides query tools, a PyTorch Geometric dataset class, and integration examples. This rich query interface---supporting searches by molecular formula, SMILES, InChI, elemental composition, and molecular fingerprint similarity---greatly facilitates targeted extraction of subsets for specific reactivity studies.

A primary strength of QM9star lies in its systematic coverage of reactive intermediates that are underrepresented in most quantum chemical databases. The inclusion of both global and local properties, particularly spin densities and atomic charges, makes it uniquely suited for machine learning studies of regioselectivity, bond-dissociation energies, and charge-dependent potential energy surfaces. However, the dataset is limited to equilibrium geometries derived from QM9; it does not sample transition states or non-equilibrium reactive pathways. The reliance on a single DFT functional (B3LYP-D3(BJ)) and a modest basis set (6-311+G(d,p)) may introduce systematic errors, especially for radicals and anions where dative correlation effects are important. Additionally, the restriction to molecules with up to nine heavy atoms and only C, N, O, F (plus H) limits applicability to larger or more diverse organic systems.

\subsection{MDCD20: Efficient Sampling of Reactive Configurations}

The MDCD20 dataset~\cite{ds50-1,ds50-2,ds50-3} addresses a complementary challenge: the generation of non-equilibrium reactive configurations---including transition states and high-energy intermediates---at a fraction of the computational cost required by traditional enumeration or AIMD sampling. It comprises approximately 1.4 million structures of neutral, radical, and charged molecules composed solely of H, C, N, and O atoms. The dataset is notably compact given its chemical diversity: 1,104,775 even-electron neutral species, 266,985 radicals, and 24,306 charged structures ($\pm 1$), with 94.8\% containing no more than 20 heavy atoms. The generation framework combines molecular dynamics (MD) and coordinate driving (CD) in an active learning loop (MD/CD-AL), seeded from a subset of Transition-1x and initial reactive structures. An ensemble of four DPA-3 neural network models was used for query-by-committee selection, with the model deviation $\sigma$ of atomic forces guiding the addition of new DFT calculations. The uncertainty threshold was gradually decreased from initial bounds of $[0.4, 3.0]~\mathrm{eV/\AA}$ to $[0.1, 0.6]~\mathrm{eV/\AA}$ over 50 active learning iterations.

Reference DFT calculations were performed with ORCA~6.0.68 at the gCP(TZ)-$\omega$B97X-D4/def2-TZVP level, using RIJCOSX acceleration and spin-unrestricted formalism with broken-symmetry initial guesses for radicals. Geometry optimizations employed Gaussian~16 for transition state verification via the Berny algorithm. The final MDCD-NN potential was trained using the DPA-3 architecture, incorporating total charge and spin multiplicity as frame-level features, with equal weighting of energy and force residuals.

The deliberate focus on reactive configurational space is MDCD20's foremost strength. Comparative SOAP descriptor analysis showed that despite being smaller in raw count than Transition-1x or ANI-1xBB, MDCD20 covers a broader range of local environments relevant to bond breaking and formation. This efficiency is crucial for training reactive MLPs that must generalize to unseen reaction pathways. The dataset has been demonstrated in automatic reaction pathway searches, transition state location, and enhanced sampling molecular dynamics.

However, several limitations should be noted. The dataset is not yet deposited in a public repository; the provided article does not list explicit download links or a license. The ADCR software used for the MD/CD workflow is proprietary, potentially hindering reproduction. The elemental scope is restricted to H, C, N, O, excluding important heteroatoms (e.g., sulfur, halogens) common in organic chemistry. Moreover, the dataset offers only energies, forces, and analytical Hessians---no detailed electronic properties such as charges or bond orders---which may limit its utility for training models that require local atomic descriptors beyond geometric information.

\subsection{Comparative Analysis and Suitability for Reactive MLP Development}

QM9star and MDCD20 exemplify two divergent philosophies for constructing reactive datasets. QM9star is a large, systematically enumerated collection of equilibrium structures for charged and radical species, rich in electronic properties and designed for prediction tasks on stable or approximately stable intermediates. MDCD20, in contrast, is a compact, intelligently sampled set of non-equilibrium configurations, optimized for capturing the transition regions of potential energy surfaces. Both are indispensable for different aspects of reactive MLP development: QM9star provides the broad coverage of stable minima across charge and spin states needed to train models that predict reactivity descriptors (e.g., ionization energies, electron affinities, radical stabilization energies), while MDCD20 supplies the genuinely reactive frames---cloak bond breaking, formation, and conformational rearrangements---that enable reactive MLPs to extrapolate beyond equilibrium.

From a practical perspective, QM9star's open database format, query tools, and deep learning integration make it immediately accessible for benchmarking and transfer learning. Its 39 properties per entry allow multi-task learning approaches. MDCD20's strength lies in its data efficiency: the active learning framework yielded a dataset that is nearly as diverse as much larger collections, potentially reducing training costs and overfitting. The use of a high-quality DFT reference (with dispersion and counterpoise corrections) ensures robust energies and forces for reactive events. However, the lack of public repository deposition and the reliance on proprietary sampling software represent barriers to widespread adoption.

When evaluating suitability for ML potential development, researchers must weigh the chemical scope, property richness, and accessibility of each dataset. For gas-phase organic reactions involving only H, C, N, and O, MDCD20 offers an excellent starting point for training a reactive deep potential, especially if the target application involves automated reaction network exploration. QM9star, on the other hand, is better suited for studying charge transfer, radical chemistry, and selectivity in small-molecule systems, and its open data infrastructure facilitates integration into existing ML workflows. Combined training on both datasets may yield a model with improved generalization across both equilibrium and transition-region geometries, though care must be taken to reconcile the different DFT levels and basis sets. Future dataset construction efforts should consider merging the systematic enumeration approach of QM9star with the active learning strategy of MDCD20 to produce datasets that are both chemically exhaustive and dynamically informative.

\section*{References}
\begin{thebibliography}{999}
\bibitem{ds1-1} Ramakrishnan R, Dral P O, Rupp M and Von Lilienfeld O A 2014 Quantum chemistry structures and properties of 134 kilo molecules Sci. Data 1 1-7. https://doi.org/10.1038/sdata.2014.22
\bibitem{ds1-52} Ruddigkeit L, Van Deursen R, Blum L C and Reymond J-L 2012 Enumeration of 166 billion organic small molecules in the chemical Universe database GDB-17 J. Chem. Inf. Model. 52 2864-75
\bibitem{ds1-53} Stewart J J 1990 MOPAC: a semiempirical molecular orbital program J. Comput. Aided Mol. Des. 4 1-103
\bibitem{ds1-54} Beck A D 1993 Density-functional thermochemistry. III. The role of exact exchange J. Chem. Phys 98 5648-6
\bibitem{ds1-55} Stephens P J, Devlin F J, Chabalowski C F and Frisch M J 1994 Ab initio calculation of vibrational absorption and circular dichroism spectra using density functional force fields J. Chem. Phys. 98 11623-7
\bibitem{ds1-56} Ditchfield R, Hehre W J and Pople J A 1971 Self-consistent molecular-orbital methods. IX. An extended Gaussian-type basis for molecular-orbital studies of organic molecules J. Chem. Phys. 54 724-8
\bibitem{ds1-57} Krishnan R, Binkley J S, Seeger R and Pople J A 1980 Self-consistent molecular orbital methods. XX. A basis set for correlated wave functions J. Chem. Phys. 72 650-4
\bibitem{ds1-58} Frisch M J et al 2009 Gaussian 09, revision d. 01 (Gaussian Inc) vol 201
\bibitem{ds1-59} Curtiss L A, Redfern P C and Raghavachari K 2007 Gaussian-4 theory using reduced order perturbation theory J. Chem. Phys. 127 124105
\bibitem{ds1-60} Ramakrishnan R, Dral P O, Rupp M and Von Lilienfeld O A 2014 Quantum chemistry structures and properties of 134 kilo molecules Figshare (available at: https://doi.org/10.6084/m9.Figshare.978904) (Accessed 7 April 2024)
\bibitem{ds10-52} Ruddigkeit L, Van Deursen R, Blum L C and Reymond J-L 2012 Enumeration of 166 billion organic small molecules in the chemical Universe database GDB-17 J. Chem. Inf. Model. 52 2864-75
\bibitem{ds10-6} Chen G et al 2024 Alchemy: a quantum chemistry dataset for benchmarking AI models (arXiv:1906.09427 [cs.LG])
\bibitem{ds10-94} Chen G, Chen P, Hsieh C-Y, Lee C-K, Liao B, Liao R, Liu W, Qiu J, Sun Q, Tang J, Zemel R and Zhang S 2024 Alchemy data set GitHub (available at: https://github.com/tencent-alchemy/Alchemy) (Accessed 26 September 2024)
\bibitem{ds11-31} Mathiasen A, Helal H, Klaser K, Balanca P, Dean J, Luschi C, Beaini D, Fitzgibbon A and Masters D 2023 Generating QM1B with PySCFIPU Advances in Neural Information Processing Systems 36 55036-55050
\bibitem{ds11-95} Fink T, Bruggesser H and Reymond J-L 2005 Virtual exploration of the small-molecule chemical Universe below 160 daltons Angew. Chem., Int. Ed. 44 1504-8
\bibitem{ds11-96} Fink T and Reymond J-L 2007 Virtual exploration of the chemical Universe up to 11 atoms of c, n, o, f: assembly of 26.4 million structures (110.9 million stereoisomers) and analysis for new ring systems, stereochemistry, physicochemical properties, compound classes and drug discovery J. Chem., Inf. Model. 47 342-53
\bibitem{ds11-97} Landrum G et al RDKit 2022(available at: https://www.rdkit.org/) (Accessed 7 July 2024)
\bibitem{ds11-98} Mathiasen A, Helal H, Klaeser K, Balanca P, Dean J, Luschi C, Beaini D, Fitzgibbon A W and Masters D 2023 QM1B dataset GitHub (available at: https://github.com/graphcore-research/qm1b-dataset) (Accessed 16 April 2024)
\bibitem{ds12-100} Najibi A and Goerigk L 2018 The nonlocal kernel in van der Waals density functionals as an additive correction: an extensive analysis with special emphasis on the B97M-V and B97M-V approaches J. Chem. Theory Comput. 14 5725-38
\bibitem{ds12-101} Eastman P, Behara P K, Dotson D L, Galvelis R, Herr J E, Horton J T, Mao Y, Chodera J D, Pritchard B P, Wang Y, Fabritiis G D and Markland T E 2022 SPICE 1.1.2 (Zenodo)(available at: https://doi.org/10.5281/zenodo.7338495) (Accessed 16 April 2024)
\bibitem{ds12-43} Eastman P et al 2023 Spice, a dataset of drug-like molecules and peptides for training machine learning potentials Sci. Data 10 11
\bibitem{ds12-44} Eastman P, Pritchard B P, Chodera J D and Markland T E 2024 Nutmeg and SPICE: models and data for biomolecular machine learning (arXiv:2406.13112v1 [physics.chem-ph])
\bibitem{ds12-92} Weigend F and Ahlrichs R 2005 Balanced basis sets of split valence, triple zeta valence and quadruple zeta valence quality for H to Rn: design and assessment of accuracy Phys. Chem. Chem. Phys. 7 3297-305
\bibitem{ds12-99} Smith D G A et al 2020 PSI4 1.4: open-source software for high-throughput quantum chemistry J. Chem. Phys. 152 184108
\bibitem{ds13-102} Kim S et al 2023 PubChem 2023 update Nucl. Acids Res. 51 1373-80
\bibitem{ds13-103} Nakata M and Shimazaki T 2017 PubChemQC database Originally accessible at https://pubchemqc.riken.jp, the dataset has been temporarily relocated to https://nakatamaho.riken.jp/pubchemqc.riken.jp/ (Accessed 17 April 2024)
\bibitem{ds13-28} Nakata M and Shimazaki T 2017 PubChemQC project: a large-scale first-principles electronic structure database for data-driven chemistry J. Chem. Inf. Model. 57 1300-8
\bibitem{ds14-102} Kim S et al 2023 PubChem 2023 update Nucl. Acids Res. 51 1373-80
\bibitem{ds14-104} Nakata M, Shimazaki T, Hashimoto M and Maeda T 2020 PubChemQC PM6 data sets Originally accessible at http://pubchemqc. riken.jp/pm6\_datasets.html, the dataset has been temporarily relocated to https://nakatamaho.riken.jp/pubchemqc.riken.jp/ pm6\_datasets.html (Accessed 18 April 2024)
\bibitem{ds14-29} Nakata M, Shimazaki T, Hashimoto M and Maeda T 2020 PubChemQC PM6: data sets of 221 million molecules with optimized molecular geometries and electronic properties J. Chem. Inf. Model. 60 5891-5899
\bibitem{ds15-105} Nakata M and Maeda T 2023 PubChemQC B3LYP/6-31G//PM6 Data Set. The PubChemQC B3LYP/6-31G//PM6 data set is freely available for download at https://nakatamaho.riken.jp/pubchemqc.riken.jp/b3lyp\_pm6\_datasets.html under a creative commons license (Accessed 23 May 2024)
\bibitem{ds15-27} Nakata M and Maeda T 2023 PubChemQC B3LYP/6-31G//PM6 data set: the electronic structures of 86 million molecules using B3LYP/6-31G calculations J. Chem. Inf. Model. 63 5734-54
\bibitem{ds15-29} Nakata M, Shimazaki T, Hashimoto M and Maeda T 2020 PubChemQC PM6: data sets of 221 million molecules with optimized molecular geometries and electronic properties J. Chem. Inf. Model. 60 5891-5899
\bibitem{ds16-106} Glavatskikh M, Leguy J, Hunault G, Cauchy T and Da Mota B 2019 PC9 dataset Figshare (available at: https://doi.org/10.6084/ m9.figshare.9033977.v1) (Accessed 16 April 2024)
\bibitem{ds16-107} Glavatskikh M, Leguy J, Hunault G, Cauchy T and Da Mota B 2019 PC9 dataset Zenodo(available at: https://doi.org/10.5281/ zenodo.3588370) (Accessed 16 April 2024)
\bibitem{ds16-26} Glavatskikh M, Leguy J, Hunault G, Cauchy T and Da Mota B 2019 Dataset's chemical diversity limits the generalizability of machine learning predictions J. Cheminf. 11 1-15
\bibitem{ds16-28} Nakata M and Shimazaki T 2017 PubChemQC project: a large-scale first-principles electronic structure database for data-driven chemistry J. Chem. Inf. Model. 57 1300-8
\bibitem{ds17-1} Ramakrishnan R, Dral P O, Rupp M and Von Lilienfeld O A 2014 Quantum chemistry structures and properties of 134 kilo molecules Sci. Data 1 1-7. https://doi.org/10.1038/sdata.2014.22
\bibitem{ds17-108} Senthil S, Chakraborty S and Ramakrishnan R 2021 Troubleshooting unstable molecules in chemical space Chem. Sci. 12 5566-73
\bibitem{ds17-109} Chai J-D and Head-Gordon M 2008 Long-range corrected hybrid density functionals with damped atom-atom dispersion corrections Phys. Chem. Chem. Phys. 10 6615-20
\bibitem{ds17-110} Kayastha P, Chakraborty S and Ramakrishnan R 2022 The bigQM7 dataset The core data, encompassing structures, ground state properties, and electronic spectra, is conveniently available for download at https://moldisgroup.github.io/bigQM7w. For deeper analysis, the NOMAD repository at https://dx.doi.org/10.17172/NOMAD/2021.09.30-1 provides the input and output files from the corresponding calculations. Additionally, a data-mining platform is available at https://moldis.tifrh.res.in/index.html
\bibitem{ds17-77} Blum L C and Reymond J-L 2009 970 million druglike small molecules for virtual screening in the chemical Universe database GDB-13 J. Am. Chem. Soc. 131 8732-3
\bibitem{ds17-85} Rappé A K, Casewit C J, Colwell K, Goddard W A III and Skiff W M 1992 UFF, a full periodic table force field for molecular mechanics and molecular dynamics simulations J. Am. Chem. Soc. 114 10024-35
\bibitem{ds17-9} Kayastha P, Chakraborty S and Ramakrishnan R 2022 The resolution-vs.-accuracy dilemma in machine learning modeling of electronic excitation spectra Digit. Discov. 1 689-702
\bibitem{ds17-92} Weigend F and Ahlrichs R 2005 Balanced basis sets of split valence, triple zeta valence and quadruple zeta valence quality for H to Rn: design and assessment of accuracy Phys. Chem. Chem. Phys. 7 3297-305
\bibitem{ds17-95} Fink T, Bruggesser H and Reymond J-L 2005 Virtual exploration of the small-molecule chemical Universe below 160 daltons Angew. Chem., Int. Ed. 44 1504-8
\bibitem{ds17-96} Fink T and Reymond J-L 2007 Virtual exploration of the chemical Universe up to 11 atoms of c, n, o, f: assembly of 26.4 million structures (110.9 million stereoisomers) and analysis for new ring systems, stereochemistry, physicochemical properties, compound classes and drug discovery J. Chem., Inf. Model. 47 342-53
\bibitem{ds18-111} Mendez D et al 2019 ChEMBL: towards direct deposition of bioassay data Nucl. Acids Res. 47 930-40
\bibitem{ds18-112} Grimme S, Bannwarth C and Shushkov P 2017 A robust and accurate tight-binding quantum chemical method for structures, vibrational frequencies and noncovalent interactions of large molecular systems parametrized for all spd-block elements (z = 1-86) J. Chem. Theory Comput. 13 1989-2009
\bibitem{ds18-113} Bannwarth C, Ehlert S and Grimme S 2019 GFN2-xTB--an accurate and broadly parametrized self-consistent tight-binding quantum chemical method with multipole electrostatics and density-dependent dispersion contributions J. Chem. Theory Comput. 15 1652-71
\bibitem{ds18-114} Isert C, Atz K, Jiménez-Luna J and Schneider G 2022 QMugs, quantum mechanical properties of drug-like molecules ETH library collection service (available at: https://doi.org/10.3929/ethz-b-000482129) (Accessed 23 May 2024)
\bibitem{ds18-41} Isert C, Atz K, Jiménez-Luna J and Schneider G 2022 QMugs, quantum mechanical properties of drug-like molecules Sci. Data 9 273
\bibitem{ds18-92} Weigend F and Ahlrichs R 2005 Balanced basis sets of split valence, triple zeta valence and quadruple zeta valence quality for H to Rn: design and assessment of accuracy Phys. Chem. Chem. Phys. 7 3297-305
\bibitem{ds19-112} Grimme S, Bannwarth C and Shushkov P 2017 A robust and accurate tight-binding quantum chemical method for structures, vibrational frequencies and noncovalent interactions of large molecular systems parametrized for all spd-block elements (z = 1-86) J. Chem. Theory Comput. 13 1989-2009
\bibitem{ds19-115} Zdrazil B et al 2024 The ChEMBL database in 2023: a drug discovery platform spanning multiple bioactivity data types and time periods Nucl. Acids Res. 52 1180-92
\bibitem{ds19-116} Folmsbee D and Hutchison G 2021 Assessing conformer energies using electronic structure and machine learning methods Int. J. Quantum Chem. 121 26381
\bibitem{ds19-117} Jurec ka P, Sponer J, Cerny J and Hobza P 2006 Benchmark database of accurate (MP2 and CCSD(T) complete basis set limit) interaction energies of small model complexes, dna base pairs and amino acid pairs Phys. Chem. Chem. Phys. 8 1985-93
\bibitem{ds19-118} Burns L A, Faver J C, Zheng Z, Marshall M S, Smith D G, Vanommeslaeghe K, MacKerell A D, Merz K M and Sherrill C D 2017 The BioFragment Database (BFDb): an open-data platform for computational chemistry analysis of noncovalent interactions J. Chem. Phys. 147 161727
\bibitem{ds19-119} Christensen A S et al 2021 OrbNet denali training data Figshare https://doi.org/10.6084/m9.Figshare.14883867 (Accessed 27 May 2024)
\bibitem{ds19-25} Christensen A S et al 2021 OrbNet Denali: A machine learning potential for biological and organic chemistry with semi-empirical cost and DFT accuracy J. Chem. Phys. 155 204103
\bibitem{ds19-7} Smith J S, Isayev O and Roitberg A E 2017 ANI-1, a data set of 20 million calculated off-equilibrium conformations for organic molecules Sci. Data 4 1-8
\bibitem{ds2-36} Kim H, Park J Y and Choi S 2019 Energy refinement and analysis of structures in the QM9 database via a highly accurate quantum chemical method Sci. Data 6 109
\bibitem{ds2-59} Curtiss L A, Redfern P C and Raghavachari K 2007 Gaussian-4 theory using reduced order perturbation theory J. Chem. Phys. 127 124105
\bibitem{ds2-61} Curtiss L A, Redfern P C and Raghavachari K 2007 Gaussian-4 theory J. Chem. Phys. 126 084108
\bibitem{ds2-62} Frisch M J et al 2016 Gaussian 16 Revision C.01 (Gaussian Inc)
\bibitem{ds2-63} Kim H, Park J Y and Choi S 2019 Highly accurate G4(MP2) benchmark on QM9 database: energy refinement and analysis of structures Figshare (available at: https://doi.org/10.6084/m9.Figshare.c.4351631.v1) (Accessed 8 April 2024)
\bibitem{ds20-120} Ernzerhof M and Scuseria G E 1999 Assessment of the Perdew-Burke-Ernzerhof exchange-correlation functional J. Chem. Phys. 110 5029-36
\bibitem{ds20-121} Chmiela S, Tkatchenko A, Sauceda H E, Poltavsky I, Schütt K T and Müller K-R 2017 Original MD17 (available at: www.sgdml. org/\#datasets) (Accessed 27 May 2024)
\bibitem{ds20-122} Christensen A S and Von Lilienfeld O A 2017 Original MD17 Figshare (available at: https://dx.doi.org/10.6084/m9.Figshare. 12672038) (Accessed 27 May 2024)
\bibitem{ds20-123} Fey M and Lenssen J E 2019 Fast graph representation learning with PyTorch geometric (arXiv:1903.02428)
\bibitem{ds20-18} Chmiela S, Tkatchenko A, Sauceda H E, Poltavsky I, Schütt K T and Müller K-R 2017 Machine learning of accurate energy-conserving molecular force fields Sci. Adv. 3 1603015
\bibitem{ds20-19} Christensen A S and Von Lilienfeld O A 2020 On the role of gradients for machine learning of molecular energies and forces Mach. Learn. Sci. Technol. 1 045018
\bibitem{ds20-20} Chmiela S, Sauceda H E, Müller K-R and Tkatchenko A 2018 Towards exact molecular dynamics simulations with machine-learned force fields Nat. Commun. 9 3887
\bibitem{ds20-92} Weigend F and Ahlrichs R 2005 Balanced basis sets of split valence, triple zeta valence and quadruple zeta valence quality for H to Rn: design and assessment of accuracy Phys. Chem. Chem. Phys. 7 3297-305
\bibitem{ds20-99} Smith D G A et al 2020 PSI4 1.4: open-source software for high-throughput quantum chemistry J. Chem. Phys. 152 184108
\bibitem{ds21-124} Mai S, Marquetand P and González L 2018 Nonadiabatic dynamics: the SHARC approach Wiley Interdiscip. Rev. Comput. Mol. Sci. 8 1370
\bibitem{ds21-125} Roos B O, Taylor P R and Sigbahn P E 1980 A complete active space SCF method (CASSCF) using a density matrix formulated super-CI approach Chem. Phys. 48 157-73
\bibitem{ds21-126} Fdez. Galván I et al 2019 OpenMolcas: From Source Code to Insight J. Chem. Theory Comput. 15 5925-64
\bibitem{ds21-127} Zhao Y and Truhlar D G 2008 The M06 suite of density functionals for main group thermochemistry, thermochemical kinetics, noncovalent interactions, excited states and transition elements: two new functionals and systematic testing of four M06-class functionals and 12 other functionals Theor. Chem. Acc. 120 215-41
\bibitem{ds21-128} Hjorth Larsen A et al 2017 The atomic simulation environment--a python library for working with atoms J. Condens. Matter Phys. 29 273002
\bibitem{ds21-129} Pengmei Z, Liu J and Shu Y Beyond MD17: the reactive xxMD dataset 2024 GitHub https://github.com/zpengmei/xxMD and Zenodo https://doi.org/10.5281/zenodo.10393859 (Accessed 27 May 2024)
\bibitem{ds21-51} Pengmei Z, Liu J and Shu Y 2024 Beyond MD17: the reactive xxMD dataset Sci. Data 11 222
\bibitem{ds22-130} Chmiela S, Vassilev-Galindo V, Unke O T, Kabylda A, Sauceda H E, Tkatchenko A and Müller K-R 2023 MD22 (available at: www. sgdml.org/\#datasets) (Accessed 27 May 2024)
\bibitem{ds22-21} Chmiela S, Vassilev-Galindo V, Unke O T, Kabylda A, Sauceda H E, Tkatchenko A and Müller K-R 2023 Accurate global machine learning force fields for molecules with hundreds of atoms Sci. Adv. 9 0873
\bibitem{ds22-75} Seifert G, Porezag D and Frauenheim T 1996 Calculations of molecules, clusters and solids with a simplified LCAO-DFT-LDA scheme Int. J. Quantum Chem. 58 185-92
\bibitem{ds22-76} Tkatchenko A, DiStasio R A Jr, Car R and Scheffler M 2012 Accurate and efficient method for many-body van der Waals interactions Phys. Rev. Lett. 108 236402
\bibitem{ds22-80} Blum V, Gehrke R, Hanke F, Havu P, Havu V, Ren X, Reuter K and Scheffler M 2009 Ab initio molecular simulations with numeric atom-centered orbitals Comput. Phys. Commun. 180 2175-96
\bibitem{ds23-131} Colomés E, Zhan Z and Oriols X 2015 Comparing Wigner, Husimi and Bohmian distributions: which one is a true probability distribution in phase space? J. Comput. Electron. 14 894-906
\bibitem{ds23-132} Zhu X, Thompson K C and Martínez T J 2019 Geodesic interpolation for reaction pathways J. Chem. Phys. 150 164103
\bibitem{ds23-133} Pinheiro M Jr, Zhang S, Dral P O and Barbatti M 2023 The WS22 database Zenodo https://doi.org/10.5281/zenodo.7032334 (Accessed 27 May 2024)
\bibitem{ds23-50} Pinheiro M Jr, Zhang S, Dral P O and Barbatti M 2023 WS22 database, wigner sampling and geometry interpolation for configurationally diverse molecular datasets Sci. Data 10 95
\bibitem{ds23-57} Krishnan R, Binkley J S, Seeger R and Pople J A 1980 Self-consistent molecular orbital methods. XX. A basis set for correlated wave functions J. Chem. Phys. 72 650-4
\bibitem{ds24-134} Werner H-J, Knowles P J, Knizia G, Manby F R and Schütz M 2012 Molpro: a general-purpose quantum chemistry program package Wiley Interdiscip. Rev. Comput. Mol. Sci. 2 242-53
\bibitem{ds24-135} Werner H-J et al 2020 The Molpro quantum chemistry package J. Chem. Phys. 152 144107
\bibitem{ds24-136} Matthews D A, Cheng L, Harding M E, Lipparini F, Stopkowicz S, Jagau T-C, Szalay P G, Gauss J and Stanton J F 2020 Coupled-cluster techniques for computational chemistry: the CFOUR program package J. Chem. Phys. 152 214108
\bibitem{ds24-137} Zhang L, Zhang S, Owens A, Yurchenko S N and Dral P O 2022 VIB5 database Figshare (available at: https://doi.org/10.6084/m9. figshare.1690328879) (Accessed 6 June 2024)
\bibitem{ds24-48} Zhang L, Zhang S, Owens A, Yurchenko S N and Dral P O 2022 VIB5 database with accurate ab initio quantum chemical molecular potential energy surfaces Sci. Data 9 84
\bibitem{ds25-138} Smith J S, Isayev O and Roitberg A E 2017 ANI-1: an extensible neural network potential with dft accuracy at force field computational cost Chem. Sci. 8 3192-203
\bibitem{ds25-139} Smith J S, Isayev O and Roitberg A E 2017 ANI-1, A data set of 20 million calculated off-equilibrium conformations for organic molecules Figshare https://doi.org/10.6084/m9.figshare.5287732.v1; GitHub https://github.com/isayev/ANI1\_dataset (Accessed 11 June 2024)
\bibitem{ds25-56} Ditchfield R, Hehre W J and Pople J A 1971 Self-consistent molecular-orbital methods. IX. An extended Gaussian-type basis for molecular-orbital studies of organic molecules J. Chem. Phys. 54 724-8
\bibitem{ds25-7} Smith J S, Isayev O and Roitberg A E 2017 ANI-1, a data set of 20 million calculated off-equilibrium conformations for organic molecules Sci. Data 4 1-8
\bibitem{ds25-78} Halgren T A 1996 Merck molecular force field. I. Basis, form, scope, parameterization and performance of MMFF94 J. Comput. Chem. 17 490-519
\bibitem{ds25-95} Fink T, Bruggesser H and Reymond J-L 2005 Virtual exploration of the small-molecule chemical Universe below 160 daltons Angew. Chem., Int. Ed. 44 1504-8
\bibitem{ds25-96} Fink T and Reymond J-L 2007 Virtual exploration of the chemical Universe up to 11 atoms of c, n, o, f: assembly of 26.4 million structures (110.9 million stereoisomers) and analysis for new ring systems, stereochemistry, physicochemical properties, compound classes and drug discovery J. Chem., Inf. Model. 47 342-53
\bibitem{ds26-111} Mendez D et al 2019 ChEMBL: towards direct deposition of bioassay data Nucl. Acids Res. 47 930-40
\bibitem{ds26-140} Smith J S, Nebgen B, Lubbers N, Isayev O and Roitberg A E 2018 Less is more: sampling chemical space with active learning J. Chem. Phys. 148 241733
\bibitem{ds26-141} Smith J S, Nebgen B T, Zubatyuk R, Lubbers N, Devereux C, Barros K, Tretiak S, Isayev O and Roitberg A E 2019 Approaching coupled cluster accuracy with a general-purpose neural network potential through transfer learning Nat. Commun. 10 2903
\bibitem{ds26-142} Smith J S, Zubatyuk R, Nebgen B, Lubbers N, Barros K, Roitberg A E, Isayev O and Tretiak S 2020 The ANI-1ccx and ANI-1x data sets, coupled-cluster and density functional theory properties for molecules Figshare (available at: https://doi.org/10.6084/m9. figshare.c.4712477; GitHub https://github.com/aiqm/ANI1x\_datasets) (Accessed 11 June 2024)
\bibitem{ds26-143} Neese F, Wennmohs F, Becker U and Riplinger C 2020 The ORCA quantum chemistry program package J. Chem. Phys. 152 224108
\bibitem{ds26-3} Smith J S, Zubatyuk R, Nebgen B, Lubbers N, Barros K, Roitberg A E, Isayev O and Tretiak S 2020 The ANI-1ccx and ANI-1x data sets, coupled-cluster and density functional theory properties for molecules Sci. Data 7 134
\bibitem{ds26-56} Ditchfield R, Hehre W J and Pople J A 1971 Self-consistent molecular-orbital methods. IX. An extended Gaussian-type basis for molecular-orbital studies of organic molecules J. Chem. Phys. 54 724-8
\bibitem{ds26-92} Weigend F and Ahlrichs R 2005 Balanced basis sets of split valence, triple zeta valence and quadruple zeta valence quality for H to Rn: design and assessment of accuracy Phys. Chem. Chem. Phys. 7 3297-305
\bibitem{ds26-95} Fink T, Bruggesser H and Reymond J-L 2005 Virtual exploration of the small-molecule chemical Universe below 160 daltons Angew. Chem., Int. Ed. 44 1504-8
\bibitem{ds26-96} Fink T and Reymond J-L 2007 Virtual exploration of the chemical Universe up to 11 atoms of c, n, o, f: assembly of 26.4 million structures (110.9 million stereoisomers) and analysis for new ring systems, stereochemistry, physicochemical properties, compound classes and drug discovery J. Chem., Inf. Model. 47 342-53
\bibitem{ds27-111} Mendez D et al 2019 ChEMBL: towards direct deposition of bioassay data Nucl. Acids Res. 47 930-40
\bibitem{ds27-144} Brauer B, Kesharwani M K, Kozuch S and Martin J M L 2016 The S66x8 benchmark for noncovalent interactions revisited: explicitly correlated ab initio methods and density functional theory Phys. Chem. Chem. Phys. 18 20905-25
\bibitem{ds27-145} Devereux C, Smith J S, Huddleston K K, Barros K, Zubatyuk R, Isayev O and Roitberg A E 2020 ANI-2 data set (Zenodo) (available at: https://doi.org/10.5281/zenodo.10108942) (Accessed 22 August 2024)
\bibitem{ds27-3} Smith J S, Zubatyuk R, Nebgen B, Lubbers N, Barros K, Roitberg A E, Isayev O and Tretiak S 2020 The ANI-1ccx and ANI-1x data sets, coupled-cluster and density functional theory properties for molecules Sci. Data 7 134
\bibitem{ds27-8} Devereux C, Smith J S, Huddleston K K, Barros K, Zubatyuk R, Isayev O and Roitberg A E 2020 Extending the applicability of the ANI deep learning molecular potential to sulfur and halogens J. Chem. Theory Comput. 16 4192-202
\bibitem{ds27-95} Fink T, Bruggesser H and Reymond J-L 2005 Virtual exploration of the small-molecule chemical Universe below 160 daltons Angew. Chem., Int. Ed. 44 1504-8
\bibitem{ds27-96} Fink T and Reymond J-L 2007 Virtual exploration of the chemical Universe up to 11 atoms of c, n, o, f: assembly of 26.4 million structures (110.9 million stereoisomers) and analysis for new ring systems, stereochemistry, physicochemical properties, compound classes and drug discovery J. Chem., Inf. Model. 47 342-53
\bibitem{ds28-140} Smith J S, Nebgen B, Lubbers N, Isayev O and Roitberg A E 2018 Less is more: sampling chemical space with active learning J. Chem. Phys. 148 241733
\bibitem{ds28-143} Neese F, Wennmohs F, Becker U and Riplinger C 2020 The ORCA quantum chemistry program package J. Chem. Phys. 152 224108
\bibitem{ds28-146} Sheppard D, Terrell R and Henkelman G 2008 Optimization methods for finding minimum energy paths J. Chem. Phys. 128 134106
\bibitem{ds28-147} Grambow C A, Pattanaik L and Green W H 2020 Reactants, products and transition states of elementary chemical reactions based on quantum chemistry Sci. Data 7 137
\bibitem{ds28-148} Chai J-D and Head-Gordon M 2008 Systematic optimization of long-range corrected hybrid density functionals J. Chem. Phys. 128 084106
\bibitem{ds28-149} Henkelman G, Uberuaga B P and Jónsson H 2000 A climbing image nudged elastic band method for finding saddle points and minimum energy paths J. Chem. Phys. 113 9901-4
\bibitem{ds28-150} Smidstrup S, Pedersen A, Stokbro K and Jónsson H 2014 Improved initial guess for minimum energy path calculations J. Chem. Phys. 140 214106
\bibitem{ds28-151} Schreiner M, Bhowmik A, Vegge T, Busk J and Winther O 2022 Transition1x-a dataset for building generalizable reactive machine learning potentials Figshare (available at: https://doi.org/10.6084/m9.figshare.19614657.v4; GitHub https://gitlab.com/ matschreiner/Transition1x) (Accessed 17 July 2024)
\bibitem{ds28-47} Schreiner M, Bhowmik A, Vegge T, Busk J and Winther O 2022 Transition1x-a dataset for building generalizable reactive machine learning potentials Sci. Data 9 779
\bibitem{ds28-56} Ditchfield R, Hehre W J and Pople J A 1971 Self-consistent molecular-orbital methods. IX. An extended Gaussian-type basis for molecular-orbital studies of organic molecules J. Chem. Phys. 54 724-8
\bibitem{ds29-152} Liang, Xu Y, Liu R and Zhu X 2019 QM-sym-database GitHub (available at: https://github.com/XI-Lab/QM-sym-database) (Accessed 2 April 2024)
\bibitem{ds29-153} Liang, Xu Y, Liu R and Zhu X 2019 QM-sym-database Figshare (available at: https://doi.org/10.6084/m9.Figshare.9638093) (Accessed 2 April 2024)
\bibitem{ds29-38} Liang J, Xu Y, Liu R and Zhu X 2019 QM-sym, a symmetrized quantum chemistry database of 135 kilo molecules Sci. Data 6 213
\bibitem{ds3-22} Nandi S, Vegge T and Bhowmik A 2023 MultiXC-QM9: large dataset of molecular and reaction energies from multi-level quantum chemical methods Sci. Data 10 783
\bibitem{ds3-63} Kim H, Park J Y and Choi S 2019 Highly accurate G4(MP2) benchmark on QM9 database: energy refinement and analysis of structures Figshare (available at: https://doi.org/10.6084/m9.Figshare.c.4351631.v1) (Accessed 8 April 2024)
\bibitem{ds3-64} Bannwarth C, Caldeweyher E, Ehlert S, Hansen A, Pracht P, Seibert J, Spicher S and Grimme S 2021 Extended tight-binding quantum chemistry methods Wiley Interdiscip. Rev. Comput. Mol. Sci. 11 1493
\bibitem{ds3-65} Te Velde G, Bickelhaupt F M, Baerends E J, Fonseca Guerra C, van Gisbergen S J, Snijders J G and Ziegler T 2001 Chemistry with ADF J. Comput. Chem. 22 931-67
\bibitem{ds3-66} Perdew J P, Burke K and Ernzerhof M 1996 Generalized gradient approximation made simple Phys. Rev. Lett. 77 3865
\bibitem{ds3-67} Nandi S, Vegge T and Bhowmik A 2023 MultiXC-QM9 Figshare (available at: https://doi.org/10.11583/DTU.c.6185986.v3) (Accessed 8 April 2024)
\bibitem{ds30-154} Liang, Xu Y, Liu R and Zhu X 2020 QM-symex-database Figshare (available at: https://doi.org/10.6084/m9.Figshare.12815276) (Accessed 2 April 2024)
\bibitem{ds30-35} Ramakrishnan R, Hartmann M, Tapavicza E and Von Lilienfeld O A 2015 Electronic spectra from TDDFT and machine learning in chemical space J. Chem. Phys. 143 084111
\bibitem{ds30-39} Liang J, Ye S, Dai T, Zha Z, Gao Y and Zhu X 2020 QM-symex, update of the QM-sym database with excited state information for 173 kilo molecules Sci. Data 7 400
\bibitem{ds30-52} Ruddigkeit L, Van Deursen R, Blum L C and Reymond J-L 2012 Enumeration of 166 billion organic small molecules in the chemical Universe database GDB-17 J. Chem. Inf. Model. 52 2864-75
\bibitem{ds31-155} Khrabrov K et al 2022 nablaDFT: large-scale conformational energy and Hamiltonian prediction benchmark and dataset Phys. Chem. Chem. Phys. 24 25853-63
\bibitem{ds31-156} Polykovskiy D et al 2020 Molecular sets (MOSES): a benchmarking platform for molecular generation models Front. Pharmacol. 11 565644
\bibitem{ds31-157} Bemis G W and Murcko M A 1996 The properties of known drugs. 1. Molecular frameworks J. Med. Chem. 39 2887-93
\bibitem{ds31-158} Degen J, Wegscheid-Gerlach C, Zaliani A and Rarey M 2008 On the art of compiling and using `drug-like' chemical fragment spaces ChemMedChem 3 1503
\bibitem{ds31-159} Barnard J M and Downs G M 1992 Clustering of chemical structures on the basis of two-dimensional similarity measures J. Chem. Inf. Comput. Sci. 32 644-9
\bibitem{ds31-23} Khrabrov K, Ber A, Tsypin A, Ushenin K, Rumiantsev E, Telepov A, Protasov D, Shenbin I, Alekseev A, Shirokikh M, Nikolenko S, Tutubalina E and Kadurin A 2024 ∇²DFT: a universal quantum chemistry dataset of drug-like molecules and a benchmark for neural network potentials Advances in Neural Information Processing Systems 37 36869-36889 https://doi.org/10.52202/079017-1162
\bibitem{ds32-12} Wahab A, Pfuderer L, Paenurk E and Gershoni-Poranne R 2022 The COMPAS project: a computational database of polycyclic aromatic systems. phase 1: cata-condensed polybenzenoid hydrocarbons J. Chem. Inf. Model. 62 3704-13
\bibitem{ds32-13} Mayo Yanes E, Chakraborty S and Gershoni-Poranne R 2024 COMPAS-2: a dataset of cata-condensed hetero-polycyclic aromatic systems Sci. Data 11 97
\bibitem{ds32-14} Wahab A and Gershoni-Poranne R 2024 COMPAS-3: a dataset of peri-condensed polybenzenoid hydrocarbons Phys. Chem. Chem. Phys. 26 15344-57
\bibitem{ds32-160} Brinkmann G, Friedrichs O, Lisken S, Peeters A and Van Cleemput N 2010 CaGe - a virtual environment for studying some special classes of plane graphs - an update MATCH Commun. Math. Comput. Chem. 63 533-52
\bibitem{ds32-161} Wahab A, Pfuderer L, Paenurk E and Gershoni-Poranne R 2022 The COMPAS project GitLab (available at: https://gitlab.com/ porannegroup/compas and https://compas.net.technion.ac.il/) (Accessed 26 September 2024)
\bibitem{ds33-15} Grambow C A, Weir H, Cunningham C N, Biancalani T and Chuang K V 2024 CREMP: conformer-rotamer ensembles of macrocyclic peptides for machine learning Sci. Data 11 859
\bibitem{ds33-162} Li J, Yanagisawa K, Sugita M, Fujie T, Ohue M and Akiyama Y 2023 CycPeptMPDB: a comprehensive database of membrane permeability of cyclic peptides J. Chem. Inf. Model. 63 2240-50
\bibitem{ds33-163} Riniker S and Landrum G A 2015 Better informed distance geometry: using what we know to improve conformation generation J. Chem. Inf. Model. 55 2562-74
\bibitem{ds33-164} Pracht P, Bohle F and Grimme S 2020 Automated exploration of the low-energy chemical space with fast quantum chemical methods Phys. Chem. Chem. Phys. 22 7169-92
\bibitem{ds33-165} Grambow C A, Weir H, Cunningham C N, Biancalani T and Chuang K V 2024 CREMP data sets Zenodo https://doi.org/10.5281/ zenodo.7931444 and https://doi.org/10.5281/zenodo.10798261; GitHub https://github.com/Genentech/cremp (Accessed 26 September 2024)
\bibitem{ds33-64} Bannwarth C, Caldeweyher E, Ehlert S, Hansen A, Pracht P, Seibert J, Spicher S and Grimme S 2021 Extended tight-binding quantum chemistry methods Wiley Interdiscip. Rev. Comput. Mol. Sci. 11 1493
\bibitem{ds33-78} Halgren T A 1996 Merck molecular force field. I. Basis, form, scope, parameterization and performance of MMFF94 J. Comput. Chem. 17 490-519
\bibitem{ds33-97} Landrum G et al RDKit 2022(available at: https://www.rdkit.org/) (Accessed 7 July 2024)
\bibitem{ds34-1} Ramakrishnan R, Dral P O, Rupp M and Von Lilienfeld O A 2014 Quantum chemistry structures and properties of 134 kilo molecules Sci. Data 1 1-7. https://doi.org/10.1038/sdata.2014.22
\bibitem{ds34-16} Axelrod S and Gomez-Bombarelli R 2022 GEOM, energy-annotated molecular conformations for property prediction and molecular generation Sci. Data 9 185
\bibitem{ds34-164} Pracht P, Bohle F and Grimme S 2020 Automated exploration of the low-energy chemical space with fast quantum chemical methods Phys. Chem. Chem. Phys. 22 7169-92
\bibitem{ds34-166} Wu Z, Ramsundar B, Feinberg E N, Gomes J, Geniesse C, Pappu A S, Leswing K and Pande V 2018 Moleculenet: a benchmark for molecular machine learning Chem. Sci. 9 513-30
\bibitem{ds34-167} Grimme S, Bohle F, Hansen A, Pracht P, Spicher S and Stahn M 2021 Efficient quantum chemical calculation of structure ensembles and free energies for nonrigid molecules J. Chem. Phys. A 125 4039-54
\bibitem{ds34-168} Axelrod S and Gomez-Bombarelli R 2022 GEOM on GitHub GitHub (available at: https://github.com/learningmatter-mit/geom) (Accessed 26 September 2024)
\bibitem{ds35-169} Groom C R, Bruno I J, Lightfoot M P and Ward S C 2016 The Cambridge structural database Struct. Sci. 72 171-9
\bibitem{ds35-170} Kneiding H, Lukin R, Lang L, Reine S, Pedersen T B, De Bin R and Balcells D 2023 Deep learning metal complex properties with natural quantum graphs Digit. Discov. 2 618-33
\bibitem{ds35-171} Kneiding H, Nova A and Balcells D 2024 Directional multiobjective optimization of metal complexes at the billion-system scale Nat. Comput. Sci. 4 263-73
\bibitem{ds35-46} Balcells D and Skjelstad B B 2020 tmQM dataset--quantum geometries and properties of 86k transition metal complexes J. Chem. Inf. Model. 60 6135-46
\bibitem{ds36-102} Kim S et al 2023 PubChem 2023 update Nucl. Acids Res. 51 1373-80
\bibitem{ds36-172} Célerse F, Wodrich M D, Vela S, Gallarati S, Fabregat R, Juraskova V and Corminboeuf C 2024 The OFF-ON database Materials cloud (available at: https://archive.materialscloud.org/record/2023.189) (Accessed 26 September 2024)
\bibitem{ds36-24} Célerse F, Wodrich M D, Vela S, Gallarati S, Fabregat R, Juraskova V and Corminboeuf C 2024 From organic fragments to photoswitchable catalysts: the OFF-ON structural repository for transferable kernel-based potentials J. Chem. Inf. Model. 64 1201-12
\bibitem{ds37-10} Schütt K T, Arbabzadah F, Chmiela S, Müller K-R and Tkatchenko A 2017 Quantum-chemical insights from deep tensor neural networks Nat. Commun. 8 13890
\bibitem{ds37-185} C7O2H10-17 and ISO17 data access page: http://quantum-machine.org/datasets/
\bibitem{ds38-17} Schütt K T, Kindermans P-J, Sauceda H E, Chmiela S, Tkatchenko A and Müller K-R 2017 SchNet: a continuous-filter convolutional neural network for modeling quantum interactions Advances in Neural Information Processing Systems pp 992-1002
\bibitem{ds38-185} C7O2H10-17 and ISO17 data access page: http://quantum-machine.org/datasets/
\bibitem{ds39-164} Pracht P, Bohle F and Grimme S 2020 Automated exploration of the low-energy chemical space with fast quantum chemical methods Phys. Chem. Chem. Phys. 22 7169-92
\bibitem{ds39-173} McKay B D, Yirik M A and Steinbeck C 2022 Surge: a fast open-source chemical graph generator J. Cheminf. 14 24
\bibitem{ds39-174} Kent P R et al 2020 QMCPACK: advances in the development, efficiency and application of auxiliary field and real-space variational and diffusion quantum Monte Carlo J. Chem. Phys. 152 174105
\bibitem{ds39-186} VQM24 data access: Zenodo https://doi.org/10.5281/zenodo.11164951
\bibitem{ds39-49} Khan D, Benali A, Kim S Y, von Rudorff G F and von Lilienfeld O A 2024 Towards comprehensive coverage of chemical space: quantum mechanical properties of 836k constitutional and conformational closed shell neutral isomers consisting of HCNOFSiPSClBr (arXiv:2405.05961)
\bibitem{ds39-97} Landrum G et al RDKit 2022(available at: https://www.rdkit.org/) (Accessed 7 July 2024)
\bibitem{ds39-99} Smith D G A et al 2020 PSI4 1.4: open-source software for high-throughput quantum chemistry J. Chem. Phys. 152 184108
\bibitem{ds4-42} Khan D, Price A J A, Ach M L and von Lilienfeld O A 2024 Adaptive hybrid density functionals (arXiv:2402.14793 [physics.chem-ph])
\bibitem{ds4-68} Sun Q 2015 Libcint: an efficient general integral library for Gaussian basis functions J. Comput. Chem. 36 1664-71
\bibitem{ds4-69} Sun Q et al 2017 PySCF: the Python-based simulations of chemistry framework WIREs Comput. Mol. Sci. 8 1340
\bibitem{ds4-70} Sun Q et al 2020 Recent developments in the PySCF program package J. Chem. Phys. 153 024109
\bibitem{ds4-71} Khan D 2024 aPBE0 GitHub (available at: https://github.com/dkhan42/aPBE0) (Accessed 26 September 2024)
\bibitem{ds4-72} Khan D and von Lilienfeld A 2024 Revised QM9 dataset ( Zenodo) (available at: https://zenodo.org/doi/10.5281/ zenodo. 10689883) (Accessed 26 September 2024)
\bibitem{ds40-175} Grimme S, Ehrlich S and Goerigk L 2011 Effect of the damping function in dispersion corrected density functional theory J. Comput. Chem. 32 1456-65
\bibitem{ds40-187} QCDGE data access: http://langroup.site/QCDGE/
\bibitem{ds40-30} Zhu Y, Li M, Xu C and Lan Z 2024 Quantum chemistry dataset with ground-and excited-state properties of 450 kilo molecules Sci. Data 11 948
\bibitem{ds41-188} QM-22 data access: GitHub https://github.com/jmbowma/QM-22
\bibitem{ds41-40} Bowman J M, Qu C, Conte R, Nandi A, Houston P L and Yu Q 2022 The MD17 datasets from the perspective of datasets for gas-phase "small" molecule potentials J. Chem. Phys. 156 240901
\bibitem{ds42-11} Vinod V and Zaspel P 2024 CheMFi: a multifidelity dataset of quantum chemical properties of diverse molecules (arXiv:2406. 14149)
\bibitem{ds42-189} CheMFi data access: GitHub https://github.com/SM4DA/CheMFi
\bibitem{ds43-190} QM9S data access: Figshare https://doi.org/10.6084/m9.figshare.24235333
\bibitem{ds43-37} Zou Z, Zhang Y, Liang L, Wei M, Leng J, Jiang J, Luo Y and Hu W 2023 A deep learning model for predicting selected organic molecular spectra Nat. Comput. Sci. 3 957-64
\bibitem{ds44-176} Pence H E and Williams A 2010 ChemSpider: an online chemical information resource J. Chem. Educ. 87 1123-4
\bibitem{ds44-177} Herr J E, Yao K, McIntyre R, Toth D W and Parkhill J 2018 Metadynamics for training neural network model chemistries: a competitive assessment J. Chem. Phys. 148 241710
\bibitem{ds44-178} Shao Y et al 2015 Advances in molecular quantum chemistry contained in the q-chem 4 program package Mol. Phys. 113 184-215
\bibitem{ds44-191} TensorMol ChemSpider former access location: https://drive.google.com/drive/folders/1IfWPs7i5kfmErIRyuhGv95dSVtNFo0e\_ (noted in the paper as no longer available)
\bibitem{ds44-45} Yao K, Herr J E, Toth D W, Mckintyre R and Parkhill J 2018 The TensorMol-0.1 model chemistry: a neural network augmented with long-range physics Chem. Sci. 9 2261-9
\bibitem{ds45-1} https://github.com/facebookresearch/fairchem
\bibitem{ds45-2} Mardirossian, N. \& Head-Gordon, M. <i>ω</i>B97M-V: A combinatorially optimized, range-separated hybrid, meta-GGA density functional with VV10 nonlocal correlation. The Journal of Chemical Physics 144 214110 (2016). https://doi.org/10.1063/1.4952647
\bibitem{ds45-3} Levine, D. S., Shuaibi, M., Taylor, M. G., Spotte-Smith, E. W. C., Hasyim, M. R., Michel, K.,. \& Wood, B. M. (2025). The Open Molecules 2025 (OMol25) Dataset, Evaluations, and Models. arXiv:2505.08762v2.
\bibitem{ds45-4} Weigend, F., \& Ahlrichs, R. (2005). Balanced basis sets of split valence, triple zeta valence and quadruple zeta valence quality for H to Rn: Design and assessment of accuracy. Physical Chemistry Chemical Physics, 7(18), 3297-3305.
\bibitem{ds45-5} Neese, F. (2012). The ORCA program system. Wiley Interdisciplinary Reviews: Computational Molecular Science, 2(1), 73-78.
\bibitem{ds45-6} https://huggingface.co/facebook/OMol25
\bibitem{ds46-1} SHNITSEL dataset on Zenodo. https://doi.org/10.5281/zenodo.15482819
\bibitem{ds46-2} XAlkeneDB dataset on Zenodo. https://doi.org/10.5281/zenodo.10890380
\bibitem{ds46-3} Curth, R., Röhrkasten, T. E., Müller, C. \& Westermayr, J. Surface Hopping Nested Instances Training Set for Excited-state Learning. Scientific Data 12 1300 (2025). https://doi.org/10.1038/s41597-025-05443-5
\bibitem{ds47-1} ANI-1xBB code repository. GitHub. https://github.com/amateurcat/ANI-1xBB
\bibitem{ds47-2} Zhang, S., Zubatyuk, R., Yang, Y., Roitberg, A. \& Isayev, O. ANI-1xBB: An ANI-Based Reactive Potential for Small Organic Molecules. Journal of Chemical Theory and Computation 21 4365-4374 (2025). https://doi.org/10.1021/acs.jctc.5c00347
\bibitem{ds47-3} Smith, J. S., Zubatyuk, R., Nebgen, B., Lubbers, N., Barros, K., Roitberg, A. E., Isayev, O. \& Tretiak, S. The ANI-1ccx and ANI-1x data sets, coupled-cluster and density functional theory properties for molecules. Scientific Data 7 134 (2020). https://doi.org/10.1038/s41597-020-0473-z
\bibitem{ds47-4} Gao, X., Ramezanghorbani, F., Isayev, O., Smith, J. S. \& Roitberg, A. E. TorchANI: A Free and Open Source PyTorch-Based Deep Learning Implementation of the ANI Neural Network Potentials. Journal of Chemical Information and Modeling 60 3408-3415 (2020). https://doi.org/10.1021/acs.jcim.0c00451
\bibitem{ds47-5} Brandenburg, J. G., Bannwarth, C., Hansen, A. \& Grimme, S. B97-3c: A revised low-cost variant of the B97-D density functional method. The Journal of Chemical Physics 148 064104 (2018). https://doi.org/10.1063/1.5012601
\bibitem{ds47-6} ANI-1xBB dataset. Carnegie Mellon University Kilthub. https://kilthub.cmu.edu/articles/dataset/ANI-1xBB\_dataset/28405316
\bibitem{ds48-1} Figshare record for QM40: https://doi.org/10.6084/m9.figshare.25993060.v1
\bibitem{ds48-2} Madushanka, A., Moura, R. T. \& Kraka, E. QM40, Realistic Quantum Mechanical Dataset for Machine Learning in Molecular Science. Scientific Data 11 1376 (2024). https://doi.org/10.1038/s41597-024-04206-y
\bibitem{ds48-3} Ramakrishnan, R., Dral, P. O., Rupp, M. \& von Lilienfeld, O. A. Quantum chemistry structures and properties of 134 kilo molecules. Scientific Data 1 140022 (2014). https://doi.org/10.1038/sdata.2014.22
\bibitem{ds48-4} Irwin, J. J., Sterling, T., Mysinger, M. M., Bolstad, E. S. \& Coleman, R. G. ZINC: A Free Tool to Discover Chemistry for Biology. Journal of Chemical Information and Modeling 52 1757-1768 (2012). https://doi.org/10.1021/ci3001277
\bibitem{ds48-5} Bannwarth, C., Ehlert, S. \& Grimme, S. GFN2-xTB—An Accurate and Broadly Parametrized Self-Consistent Tight-Binding Quantum Chemical Method with Multipole Electrostatics and Density-Dependent Dispersion Contributions. Journal of Chemical Theory and Computation 15 1652-1671 (2019). https://doi.org/10.1021/acs.jctc.8b01176
\bibitem{ds49-1} Tang, M., Zhu, T., Zhang, S. \& Hong, X. QM9star, two Million DFT-computed Equilibrium Structures for Ions and Radicals with Atomic Information. figshare https://doi.org/10.6084/m9.figshare.27002905 (2024).
\bibitem{ds49-2} Tang, M., Zhu, T., Zhang, S. \& Hong, X. QM9star query code repository. GitHub https://github.com/gentle1999/qm9star\_query (2024).
\bibitem{ds49-3} Tang, M. J., Zhu, T. C., Zhang, S. Q. \& Hong, X. QM9star, two Million DFT-computed Equilibrium Structures for Ions and Radicals with Atomic Information. Scientific Data 11 1158 (2024). https://doi.org/10.1038/s41597-024-03933-6
\bibitem{ds49-4} Ramakrishnan, R., Dral, P. O., Rupp, M. \& von Lilienfeld, O. A. Quantum chemistry structures and properties of 134 kilo molecules. Scientific Data 1 140022 (2014). https://doi.org/10.1038/sdata.2014.22
\bibitem{ds5-5} Meza-González B, Ramírez-Palma D I, Carpio-Martínez P, Vázquez-Cuevas D, Martínez-Mayorga K and Cortés-Guzmán F 2024 Quantum topological atomic properties of 44k molecules Sci. Data 11 945
\bibitem{ds5-60} Ramakrishnan R, Dral P O, Rupp M and Von Lilienfeld O A 2014 Quantum chemistry structures and properties of 134 kilo molecules Figshare (available at: https://doi.org/10.6084/m9.Figshare.978904) (Accessed 7 April 2024)
\bibitem{ds5-73} Keith T A AIMAll (version 19.10.12). TK Gristmill Software, Overland Park KS, USA (2019(aim.tkgristmill.com))
\bibitem{ds5-74} Meza-González B, Ramírez-Palma D I, Carpio-Martínez P, Vázquez-Cuevas D, Martínez-Mayorga K and Cortés-Guzmán F 2024 AIMEl-DB data set at Zenodo (Zenodo)(available at: https://zenodo.org/records/11406726) (Accessed 26 September 2024)
\bibitem{ds50-1} Li, G., Ling, H., Su, C., Liu, Z., Wang, G., Yang, M. \& Li, S. A Data-Efficient Reactive Machine Learning Potential to Accelerate Automated Exploration of Complex Reaction Networks. CCS Chemistry 1-16 (2026). https://doi.org/10.31635/ccschem.026.202607339
\bibitem{ds50-2} Smith, J. S., Zubatyuk, R., Nebgen, B., Lubbers, N., Barros, K., Roitberg, A. E., Isayev, O. \& Tretiak, S. The ANI-1ccx and ANI-1x data sets, coupled-cluster and density functional theory properties for molecules. Scientific Data 7 134 (2020). https://doi.org/10.1038/s41597-020-0473-z
\bibitem{ds50-3} Neese, F. The ORCA program system. WIREs Computational Molecular Science 2 73-78 (2011). https://doi.org/10.1002/wcms.81
\bibitem{ds6-33} Hoja J, Medrano Sandonas L, Ernst B G, Vázquez-Mayagoitia A, DiStasio R A Jr and Tkatchenko A 2021 QM7-X, a comprehensive dataset of quantum-mechanical properties spanning the chemical space of small organic molecules Sci. Data 8 43
\bibitem{ds6-75} Seifert G, Porezag D and Frauenheim T 1996 Calculations of molecules, clusters and solids with a simplified LCAO-DFT-LDA scheme Int. J. Quantum Chem. 58 185-92
\bibitem{ds6-76} Tkatchenko A, DiStasio R A Jr, Car R and Scheffler M 2012 Accurate and efficient method for many-body van der Waals interactions Phys. Rev. Lett. 108 236402
\bibitem{ds6-77} Blum L C and Reymond J-L 2009 970 million druglike small molecules for virtual screening in the chemical Universe database GDB-13 J. Am. Chem. Soc. 131 8732-3
\bibitem{ds6-78} Halgren T A 1996 Merck molecular force field. I. Basis, form, scope, parameterization and performance of MMFF94 J. Comput. Chem. 17 490-519
\bibitem{ds6-79} O'Boyle N M, Vandermeersch T, Flynn C J, Maguire A R and Hutchison G R 2011 Confab-systematic generation of diverse low-energy conformers J. Cheminf. 3 1-9
\bibitem{ds6-80} Blum V, Gehrke R, Hanke F, Havu P, Havu V, Ren X, Reuter K and Scheffler M 2009 Ab initio molecular simulations with numeric atom-centered orbitals Comput. Phys. Commun. 180 2175-96
\bibitem{ds6-81} Hoja J, Medrano Sandonas L, Ernst B G, Vázquez-Mayagoitia A, DiStasio R A Jr and Tkatchenko A 2020 QM7-X, a comprehensive dataset of quantum-mechanical properties spanning the chemical space of small organic molecules Zenodo https://doi.org/10.5281/zenodo.4288677 (Accessed 9 April 2024)
\bibitem{ds7-32} Rupp M, Tkatchenko A, Müller K-R and Von Lilienfeld O A 2012 Fast and accurate modeling of molecular atomization energies with machine learning Phys. Rev. Lett. 108 058301
\bibitem{ds7-77} Blum L C and Reymond J-L 2009 970 million druglike small molecules for virtual screening in the chemical Universe database GDB-13 J. Am. Chem. Soc. 131 8732-3
\bibitem{ds7-82} Guha R, Howard M T, Hutchison G R, Murray-Rust P, Rzepa H, Steinbeck C, Wegner J and Willighagen E L 2006 The Blue Obelisk--interoperability in chemical informatics J. Chem. Inf. Model. 46 991-8
\bibitem{ds7-83} Rupp M, Tkatchenko A, Müller K-R and Von Lilienfeld O A 2012 QM7 dataset(available at: http://quantum-machine.org/ datasets/\#qm7) (Accessed 9 April 2024)
\bibitem{ds8-34} Montavon G, Rupp M, Gobre V, Vázquez-Mayagoitia A, Hansen K, Tkatchenko A, Müller K-R and Von Lilienfeld O A 2013 Machine learning of molecular electronic properties in chemical compound space New J. Phys. 15 095003
\bibitem{ds8-53} Stewart J J 1990 MOPAC: a semiempirical molecular orbital program J. Comput. Aided Mol. Des. 4 1-103
\bibitem{ds8-75} Seifert G, Porezag D and Frauenheim T 1996 Calculations of molecules, clusters and solids with a simplified LCAO-DFT-LDA scheme Int. J. Quantum Chem. 58 185-92
\bibitem{ds8-76} Tkatchenko A, DiStasio R A Jr, Car R and Scheffler M 2012 Accurate and efficient method for many-body van der Waals interactions Phys. Rev. Lett. 108 236402
\bibitem{ds8-84} Weininger D 1988 SMILES, a chemical language and information system. 1. Introduction to methodology and encoding rules J. Chem. Inf. Comput. Sci. 28 31-36
\bibitem{ds8-85} Rappé A K, Casewit C J, Colwell K, Goddard W A III and Skiff W M 1992 UFF, a full periodic table force field for molecular mechanics and molecular dynamics simulations J. Am. Chem. Soc. 114 10024-35
\bibitem{ds8-86} Hedin L 1965 New method for calculating the one-particle green's function with application to the electron-gas problem Phys. Rev. 139 796
\bibitem{ds8-87} Neese F 2022 Software update: The ORCA program system--Version 5.0 Wiley Interdiscip. Rev. Comput. Mol. Sci. 12 e1606
\bibitem{ds8-88} Montavon G, Rupp M, Gobre V, Vázquez-Mayagoitia A, Hansen K, Tkatchenko A, Müller K-R and Von Lilienfeld O A 2013 QM7b dataset(available at: http://stacks.iop.org/NJP/15/095003/mmedia) (Accessed 15 April 2024)
\bibitem{ds9-35} Ramakrishnan R, Hartmann M, Tapavicza E and Von Lilienfeld O A 2015 Electronic spectra from TDDFT and machine learning in chemical space J. Chem. Phys. 143 084111
\bibitem{ds9-89} Furche F, Ahlrichs R, Hüttig C, Klopper W, Sierka M and Weigend F 2014 Turbomole Wiley Interdiscip. Rev. Comput. Mol. Sci. 4 91-100
\bibitem{ds9-90} Furche F and Ahlrichs R 2002 Adiabatic time-dependent density functional methods for excited state properties J. Chem. Phys. 117 7433-47
\bibitem{ds9-91} Perdew J P, Ernzerhof M and Burke K 1996 Rationale for mixing exact exchange with density functional approximations J. Chem. Phys. 105 9982-5
\bibitem{ds9-92} Weigend F and Ahlrichs R 2005 Balanced basis sets of split valence, triple zeta valence and quadruple zeta valence quality for H to Rn: design and assessment of accuracy Phys. Chem. Chem. Phys. 7 3297-305
\bibitem{ds9-93} Hüttig C and Weigend F 2000 CC2 excitation energy calculations on large molecules using the resolution of the identity approximation J. Chem. Phys. 113 5154-61
\end{thebibliography}

\end{document}
